\chapter{Relative Mealy Program Categories and Dynamic Span Logic}
\label{ch:relative-mealy-program-categories}

\section{Overview}
\label{sec:relative-mealy-overview}

This chapter develops a span-based, non-iterative dynamic logic for Mealy
morphisms and Mealy simulations. We use the span conventions, composition
order, internal-category conventions, Mealy morphisms, output-comparison
monad, input-processing monad, and Mealy simulations fixed in
Chapter~\ref{ch:spans-internal-categories-mealy}.

Unless otherwise stated, \(\mathcal E\) is finitely complete.

The stateful morphism layer is close in spirit to Paré's Mealy morphisms of
enriched categories, which generalize classical Mealy machines
\cite{PareMealyMorphisms}. The logical layer follows the dynamic-logic
tradition of interpreting programs by diamond, box, test, and choice
modalities, but this chapter isolates the non-iterative fragment
\cite{HarelKozenTiurynDynamicLogic}. Categorically, the predicate-transformer
semantics uses the standard subobject-fibration reading of regular, coherent,
and first-order Heyting logic
\cite{Jacobs1999CategoricalLogic}; its regular fragment is closely related to
the relational calculus of regular categories
\cite{FongSpivakRegularRelational}.

As in Chapter~\ref{ch:spans-internal-categories-mealy}, all displayed
generalized-element equations denote the corresponding morphism equations out
of the appropriate pullback object.

Throughout this chapter we use the composition convention
\[
m_A(a,b)=b\circ a
\]
for internally composable arrows
\[
a:u\to v,
\qquad
b:v\to w.
\]
Similarly,
\[
m_B(\beta,\gamma)=\gamma\circ\beta.
\]
Thus right actions execute target-to-source: if \(p(x)=w\), then \(b\) acts
first and \(a\) acts second.

Let
\[
\Phi=(P,\delta_P,\omega_P):\mathbb A\rightsquigarrow\mathbb B
\]
be a Mealy morphism with carrier span
\[
\begin{tikzcd}
& P \arrow[dl,"p_P"'] \arrow[dr,"q_P"] & \\
A_0 && B_0
\end{tikzcd}
.
\]
Let
\[
Y=(Y\xrightarrow{r}B_0,\sigma)
\]
be a downstream right \(\mathbb B\)-action. The coupled state object is the
pullback
\[
\begin{tikzcd}
S_{\Phi,Y}
\arrow[r,"\pi_Y"]
\arrow[d,"\pi_P"']
&
Y
\arrow[d,"r"]
\\
P
\arrow[r,"q_P"']
&
B_0
\arrow[ul,phantom,very near start,"\lrcorner"]
\end{tikzcd}
.
\]
Thus
\[
S_{\Phi,Y}:=P\times_{q_P,B_0,r}Y.
\]
The action of an input arrow \(a:u\to v\) on a coupled state
\((x,y)\), with \(p_P(x)=v\), is given by
\[
a\cdot(x,y)
=
\bigl(\delta_P(a,x),\sigma(\omega_P(a,x),y)\bigr).
\]
The result lies over \(u\), and therefore defines a right
\(\mathbb A\)-action
\[
\Phi^\ast Y.
\]
The relative program category is the corresponding action category
\[
\mathsf{Prog}_{\Phi,Y}
:=
\int_{\mathbb A}\Phi^\ast Y.
\]

A Mealy simulation
\[
\Theta:\Phi\Rightarrow\Psi
\]
from
\[
\Phi=(P,\delta_P,\omega_P)
\qquad\text{to}\qquad
\Psi=(Q,\delta_Q,\omega_Q)
\]
is represented by an output-comparison Kleisli cell
\[
Q\longrightarrow B\circ P.
\]
Here \(B\) denotes the underlying endospan of the internal category
\(\mathbb B\). In components this is a pair
\[
\theta:Q\to P,
\qquad
\alpha:Q\to B_1,
\]
displayed by
\[
\begin{tikzcd}[column sep=large,row sep=large]
&& B_1
  \arrow[dd,bend left=18,"s_B"]
  \arrow[dd,bend right=18,"t_B"']
&
\\
Q
  \arrow[rru,"\alpha"]
  \arrow[r,"\theta"']
  \arrow[drr,"q_Q"']
  \arrow[ddrr,bend right=16,"p_Q"']
&
P
  \arrow[dr,"q_P"]
  \arrow[rdd,bend left=16,"p_P"]
&&
\\
&& B_0 &
\\
&& A_0 &
\end{tikzcd}
\]
with equations
\[
p_P\theta=p_Q,
\qquad
s_B\alpha=q_P\theta,
\qquad
t_B\alpha=q_Q.
\]
Thus each state \(y\in Q\) is compared with the state
\(\theta y\in P\) over the same input interface, and the output interfaces
are compared by an arrow
\[
q_P(\theta y)\xrightarrow{\alpha_y}q_Q(y)
\]
of \(\mathbb B\).

A downstream action \(Y\) evaluates this comparison arrow as an actual
transition in the action category \(\int_{\mathbb B}Y\). Hence
\(\Theta\) induces a map
\[
h_{\Theta,Y}:S_{\Psi,Y}\to S_{\Phi,Y}
\]
by
\[
h_{\Theta,Y}(y,z)
=
\bigl(\theta y,\sigma(\alpha_y,z)\bigr).
\]
The simulation law is precisely the assertion that this map is
\(\mathbb A\)-equivariant. Consequently it induces an internal functor
\[
H_{\Theta,Y}:
\mathsf{Prog}_{\Psi,Y}
\longrightarrow
\mathsf{Prog}_{\Phi,Y}.
\]
The relative program construction is therefore contravariant in Mealy
simulations.

The guiding construction is:
\[
\boxed{
\text{a downstream action evaluates output-comparison Kleisli cells as
state maps.}
}
\]

After the program category has been formed, dynamic span logic is obtained by
applying span predicate transformers to endospans on the coupled state object
\(S_{\Phi,Y}\). In a regular category these transformers give diamond
modalities. In a first-order Heyting category they also give box modalities.
This chapter treats the one-step and finite-choice fragment of relative
dynamic span logic: primitive execution spans, tests, restricted primitive
subprograms, diamonds, boxes, and stability under downstream change and Mealy
simulation.

\section{Right actions and action categories}
\label{sec:relative-right-actions-action-categories}

Throughout this section let
\[
\mathbb A=(A_0,A_1,s_A,t_A,e_A,m_A)
\]
be an internal category in \(\mathcal E\).

The use of actions as Eilenberg--Moore algebras over a slice, and the
corresponding subobject-fibration logic used later in the chapter, follows the
standard categorical-logic perspective on indexed predicates and their
quantifiers \cite{Jacobs1999CategoricalLogic}.

\subsection{Right internal actions}
\label{subsec:relative-right-actions}

\begin{definition}[Right internal action]
\label{def:relative-right-action}
A \textbf{right action} of \(\mathbb A\) consists of a morphism
\[
p:X\to A_0
\]
together with a map
\[
\rho:
A_1\times_{t_A,A_0,p}X\to X
\]
such that
\[
p\rho=s_A\pi_1.
\]
The typing square for the action domain is
\[
\begin{tikzcd}
A_1\times_{t_A,A_0,p}X
\arrow[r,"\pi_X"]
\arrow[d,"\pi_1"']
&
X
\arrow[d,"p"]
\\
A_1
\arrow[r,"t_A"']
&
A_0
\arrow[ul,phantom,very near start,"\lrcorner"]
\end{tikzcd}
.
\]
The action map fits into
\[
\begin{tikzcd}
A_1\times_{t_A,A_0,p}X
\arrow[rr,"\rho"]
\arrow[dr,"s_A\pi_1"']
&&
X
\arrow[dl,"p"]
\\
& A_0 &
\end{tikzcd}
\]
and satisfies the unit and multiplication laws
\[
\rho(e_A(px),x)=x
\]
and
\[
\rho(b\circ a,x)=\rho(a,\rho(b,x))
\]
for internally composable arrows
\[
a:u\to v,
\qquad
b:v\to w,
\qquad
p(x)=w.
\]
Equivalently, using the multiplication map,
\[
\rho(m_A(a,b),x)=\rho(a,\rho(b,x)).
\]
\end{definition}

\begin{definition}[Equivariant map]
\label{def:relative-equivariant-map}
Let
\[
(X\xrightarrow{p}A_0,\rho_X)
\qquad\text{and}\qquad
(Z\xrightarrow{q}A_0,\rho_Z)
\]
be right \(\mathbb A\)-actions. An \(\mathbb A\)-equivariant map
\[
f:X\to Z
\]
is a morphism in \(\mathcal E\) satisfying
\[
qf=p
\]
and
\[
f(\rho_X(a,x))=\rho_Z(a,f(x)).
\]
Equivalently, the diagrams
\[
\begin{tikzcd}
X \arrow[rr,"f"] \arrow[dr,"p"'] && Z \arrow[dl,"q"] \\
& A_0 &
\end{tikzcd}
\]
and
\[
\begin{tikzcd}
A_1\times_{t_A,A_0,p}X
\arrow[r,"1\times f"]
\arrow[d,"\rho_X"']
&
A_1\times_{t_A,A_0,q}Z
\arrow[d,"\rho_Z"]
\\
X
\arrow[r,"f"']
&
Z
\end{tikzcd}
\]
commute.
\end{definition}

\begin{definition}[Category of right actions]
\label{def:relative-action-category-external}
The \textbf{category of right \(\mathbb A\)-actions}, denoted
\[
\mathbf{Act}(\mathbb A),
\]
has right internal \(\mathbb A\)-actions as objects and equivariant maps as
morphisms.
\end{definition}

\subsection{Actions as Eilenberg--Moore algebras}
\label{subsec:relative-actions-em}

\begin{proposition}[Actions as Eilenberg--Moore algebras]
\label{prop:relative-actions-em}
The internal category \(\mathbb A\), regarded as a monad in
\(\mathbf{Span}(\mathcal E)\), induces a monad
\[
T_{\mathbb A}:\mathcal E/A_0\to\mathcal E/A_0
\]
defined by
\[
T_{\mathbb A}(X\xrightarrow{p}A_0)
=
\left(
A_1\times_{t_A,A_0,p}X
\xrightarrow{s_A\pi_1}
A_0
\right).
\]
The category \(\mathbf{Act}(\mathbb A)\) is the Eilenberg--Moore category of
\(T_{\mathbb A}\).
\end{proposition}

\begin{proof}
\leavevmode
\begin{enumerate}
\item \textbf{Identify the endofunctor.}
For an object \(p:X\to A_0\) of \(\mathcal E/A_0\), composing the span
\[
A_0\xleftarrow{s_A}A_1\xrightarrow{t_A}A_0
\]
with \(p\) gives the pullback
\[
\begin{tikzcd}
A_1\times_{t_A,A_0,p}X
\arrow[r,"\pi_X"]
\arrow[d,"\pi_1"']
&
X
\arrow[d,"p"]
\\
A_1
\arrow[r,"t_A"']
&
A_0
\arrow[ul,phantom,very near start,"\lrcorner"]
\end{tikzcd}
,
\]
equipped with the structure map \(s_A\pi_1\) to \(A_0\).

\item \textbf{Unpack an algebra.}
A \(T_{\mathbb A}\)-algebra structure on \(p:X\to A_0\) is a map
\[
A_1\times_{t_A,A_0,p}X\to X
\]
over \(A_0\). This is precisely a map \(\rho\) satisfying
\[
p\rho=s_A\pi_1.
\]

\item \textbf{Compare the laws.}
The unit of the monad is induced by \(e_A:A_0\to A_1\), and the multiplication
is induced by \(m_A:A_1\times_{A_0}A_1\to A_1\). The Eilenberg--Moore unit
and multiplication equations are exactly
\[
\rho(e_A(px),x)=x
\]
and
\[
\rho(b\circ a,x)=\rho(a,\rho(b,x)).
\]

\item \textbf{Compare the morphisms.}
A morphism of Eilenberg--Moore algebras is a map over \(A_0\) commuting with
the algebra structures. This is exactly an equivariant map.
\end{enumerate}
\end{proof}

\subsection{Action categories}
\label{subsec:relative-action-categories}

Let
\[
X=(X\xrightarrow{p}A_0,\rho)
\]
be a right \(\mathbb A\)-action.

\begin{definition}[Action category]
\label{def:relative-action-category}
The \textbf{action category} of \(X\), denoted
\[
\int_{\mathbb A}X,
\]
is the internal category with
\[
\left(\int_{\mathbb A}X\right)_0=X,
\qquad
\left(\int_{\mathbb A}X\right)_1
=
A_1\times_{t_A,A_0,p}X.
\]
The source and target maps are
\[
s(a,x)=x,
\qquad
t(a,x)=\rho(a,x).
\]
They form the span
\[
\begin{tikzcd}
& A_1\times_{t_A,A_0,p}X
\arrow[dl,"\pi_X"']
\arrow[dr,"\rho"]
&
\\
X && X
\end{tikzcd}
.
\]
The identity map is
\[
e(x)=(e_A(px),x).
\]
Composition is given by
\[
\bigl((b,x),(a,\rho(b,x))\bigr)
\longmapsto
(m_A(a,b),x).
\]
Equivalently, the composite of
\[
x\xrightarrow{\,b\,}\rho(b,x)
\xrightarrow{\,a\,}\rho(a,\rho(b,x))
\]
is labelled by
\[
b\circ a=m_A(a,b).
\]
\end{definition}

\begin{theorem}[Action categories]
\label{thm:relative-action-category-internal}
For every right \(\mathbb A\)-action \(X\), the data of
Definition~\ref{def:relative-action-category} define an internal category
\[
\int_{\mathbb A}X.
\]
Equivalently, the span
\[
X
\xleftarrow{\pi_X}
A_1\times_{t_A,A_0,p}X
\xrightarrow{\rho}
X
\]
is a monoid object in
\[
\operatorname{EndSpan}_{\mathcal E}(X).
\]
\end{theorem}

\begin{proof}
\leavevmode
\begin{enumerate}
\item \textbf{Check the target map.}
The equation \(p\rho=s_A\pi_1\) says that the target of an action arrow is
typed over the source of its input arrow. If
\[
a:u\to v,
\qquad
p(x)=v,
\]
then \(\rho(a,x)\) lies over \(u\).

\item \textbf{Check identities.}
The identity arrow at \(x\) is
\[
(e_A(px),x).
\]
It lies in the pullback because \(t_Ae_Ap=p\). The action unit law gives
\[
\rho(e_A(px),x)=x.
\]

\item \textbf{Describe composable pairs.}
A composable pair in \(\int_{\mathbb A}X\) has the form
\[
\bigl((b,x),(a,\rho(b,x))\bigr)
\]
with
\[
a:u\to v,
\qquad
b:v\to w,
\qquad
p(x)=w.
\]
The corresponding pullback diagram is
\[
\begin{tikzcd}
\left(A_1\times_{A_0}X\right)
\times_X
\left(A_1\times_{A_0}X\right)
\arrow[r]
\arrow[d]
&
A_1\times_{A_0}X
\arrow[d,"\pi_X"]
\\
A_1\times_{A_0}X
\arrow[r,"\rho"']
&
X
\arrow[ul,phantom,very near start,"\lrcorner"]
\end{tikzcd}
.
\]

\item \textbf{Check composition.}
The composite
\[
(m_A(a,b),x)
\]
is defined because
\[
t_A(m_A(a,b))=t_A(b)=p(x).
\]
Its target is
\[
\rho(m_A(a,b),x)=\rho(a,\rho(b,x))
\]
by the action multiplication law.

\item \textbf{Verify the laws.}
The identity and associativity laws follow from the unit and associativity
laws of \(\mathbb A\) together with the action laws.
\end{enumerate}
\end{proof}

\subsection{Input-discrete functors}
\label{subsec:relative-input-discrete-functors}

The action category has a canonical internal functor
\[
\pi_X:\int_{\mathbb A}X\to\mathbb A^{\mathrm{op}}
\]
whose object part is
\[
p:X\to A_0
\]
and whose arrow part is the projection
\[
\pi_1:A_1\times_{t_A,A_0,p}X\to A_1.
\]
The source square of this functor is the defining pullback
\[
\begin{tikzcd}
A_1\times_{t_A,A_0,p}X
\arrow[r,"\pi_1"]
\arrow[d,"\pi_X"']
&
A_1
\arrow[d,"t_A"]
\\
X
\arrow[r,"p"']
&
A_0
\arrow[ul,phantom,very near start,"\lrcorner"]
\end{tikzcd}
.
\]
The use of \(t_A\) rather than \(s_A\) is because the codomain is
\(\mathbb A^{\mathrm{op}}\).

\begin{definition}[Input-discrete functor]
\label{def:relative-input-discrete}
Let
\[
J:\mathbb X\to\mathbb A^{\mathrm{op}}
\]
be an internal functor with object map
\[
J_0:X_0\to A_0.
\]
We say that \(J\) is \textbf{input-discrete} if the square
\[
\begin{tikzcd}
X_1 \arrow[r,"J_1"] \arrow[d,"s_X"']
&
A_1 \arrow[d,"t_A"]
\\
X_0 \arrow[r,"J_0"']
&
A_0
\end{tikzcd}
\]
is a pullback.
\end{definition}

\begin{theorem}[Actions as input-discrete functors]
\label{thm:relative-actions-input-discrete}
To give a right \(\mathbb A\)-action
\[
(X\xrightarrow{p}A_0,\rho)
\]
is equivalently to give an internal category \(\mathbb X\) with object of
objects \(X\), together with an input-discrete internal functor
\[
J:\mathbb X\to\mathbb A^{\mathrm{op}}
\]
whose object map is \(p\).
\end{theorem}

\begin{proof}
\leavevmode
\begin{enumerate}
\item \textbf{Construct the functor from an action.}
Given a right action \(X\), take
\[
\mathbb X=\int_{\mathbb A}X.
\]
Its arrow object is
\[
A_1\times_{t_A,A_0,p}X,
\]
and the functor to \(\mathbb A^{\mathrm{op}}\) is given by the projections
\[
p:X\to A_0,
\qquad
\pi_1:A_1\times_{t_A,A_0,p}X\to A_1.
\]
The input-discrete square is the pullback defining the arrow object.

\item \textbf{Recover an action from an input-discrete functor.}
Conversely, suppose
\[
J:\mathbb X\to\mathbb A^{\mathrm{op}}
\]
is input-discrete with object map \(p:X\to A_0\). The pullback condition gives
a canonical isomorphism
\[
X_1\cong A_1\times_{t_A,A_0,p}X.
\]
Transporting the target map of \(\mathbb X\) across this isomorphism gives
\[
\rho:A_1\times_{t_A,A_0,p}X\to X.
\]

\item \textbf{Check the typing.}
The target compatibility of \(J\) gives
\[
p\rho=s_A\pi_1.
\]
This is the required action typing equation.

\item \textbf{Check the laws.}
Under the pullback identification, the identity map of \(\mathbb X\) is
\[
x\longmapsto(e_A(px),x)
\]
because \(J\) preserves identities. Hence the unit law for \(\rho\) follows
from the identity law in \(\mathbb X\). Composition preservation by \(J\)
forces the transported composition map to be
\[
\bigl((b,x),(a,\rho(b,x))\bigr)
\longmapsto
(m_A(a,b),x).
\]
Associativity in \(\mathbb X\) gives
\[
\rho(m_A(a,b),x)=\rho(a,\rho(b,x)).
\]

\item \textbf{Compare the two constructions.}
The two assignments are inverse up to the canonical pullback isomorphism.
\end{enumerate}
\end{proof}

\section{Downstream evaluation of output-comparison cells}
\label{sec:relative-kleisli-evaluation}

The construction in this section treats an output-comparison Kleisli cell as a
context-dependent comparison which becomes an actual map only after a
downstream action has been chosen.

Fix internal categories
\[
\mathbb A=(A_0,A_1,s_A,t_A,e_A,m_A),
\qquad
\mathbb B=(B_0,B_1,s_B,t_B,e_B,m_B).
\]
Let
\[
\mathcal C_{A,B}:=\mathbf{Span}(\mathcal E)(A_0,B_0)
\]
be the local span category. The output-comparison monad is
\[
T_B:=B\circ -:\mathcal C_{A,B}\to\mathcal C_{A,B}.
\]
For a span
\[
P=\left(A_0\xleftarrow{p_P}P\xrightarrow{q_P}B_0\right),
\]
the object \(T_B(P)=B\circ P\) has apex
\[
P\times_{q_P,B_0,s_B}B_1
\]
and fits into
\[
\begin{tikzcd}
P\times_{q_P,B_0,s_B}B_1
\arrow[r,"\pi_{B_1}"]
\arrow[d,"\pi_P"']
&
B_1
\arrow[d,"s_B"]
\\
P
\arrow[r,"q_P"']
&
B_0
\arrow[ul,phantom,very near start,"\lrcorner"]
\end{tikzcd}
.
\]

Let
\[
Y=(Y\xrightarrow{r}B_0,\sigma)
\]
be a right \(\mathbb B\)-action.

\subsection{Evaluation on carriers}
\label{subsec:relative-evaluation-carriers}

\begin{definition}[Downstream evaluation on carriers]
\label{def:relative-downstream-evaluation-carriers}
For a span
\[
P=
\left(
A_0\xleftarrow{p_P}P\xrightarrow{q_P}B_0
\right)
\]
define
\[
\operatorname{Ev}_Y(P)
:=
P\times_{q_P,B_0,r}Y
\xrightarrow{p_P\pi_P}
A_0.
\]
The underlying pullback is
\[
\begin{tikzcd}
\operatorname{Ev}_Y(P)
\arrow[r,"\pi_Y"]
\arrow[d,"\pi_P"']
&
Y
\arrow[d,"r"]
\\
P
\arrow[r,"q_P"']
&
B_0
\arrow[ul,phantom,very near start,"\lrcorner"]
\end{tikzcd}
.
\]
\end{definition}

\subsection{Evaluation of Kleisli comparisons}
\label{subsec:relative-evaluation-kleisli-comparisons}

\begin{definition}[Evaluation of a Kleisli comparison]
\label{def:relative-kleisli-comparison-evaluation}
Let
\[
\Gamma:Q\to T_B(P)=B\circ P
\]
be a Kleisli arrow in \(\mathrm{Kl}(T_B)\). In components, write
\[
\Gamma(u)=(\gamma u,\lambda_u),
\]
where
\[
\gamma:Q\to P,
\qquad
\lambda:Q\to B_1,
\]
and
\[
p_P\gamma=p_Q,
\qquad
s_B\lambda=q_P\gamma,
\qquad
t_B\lambda=q_Q.
\]
The component data are displayed by
\[
\begin{tikzcd}[column sep=large,row sep=large]
&& B_1
  \arrow[dd,bend left=18,"s_B"]
  \arrow[dd,bend right=18,"t_B"']
&
\\
Q
  \arrow[rru,"\lambda"]
  \arrow[r,"\gamma"']
  \arrow[drr,"q_Q"']
  \arrow[ddrr,bend right=16,"p_Q"']
&
P
  \arrow[dr,"q_P"]
  \arrow[rdd,bend left=16,"p_P"]
&&
\\
&& B_0 &
\\
&& A_0 &
\end{tikzcd}
.
\]
Define
\[
\operatorname{Ev}_Y(\Gamma):
Q\times_{q_Q,B_0,r}Y
\longrightarrow
P\times_{q_P,B_0,r}Y
\]
by
\[
\operatorname{Ev}_Y(\Gamma)(u,z)
=
\bigl(\gamma u,\sigma(\lambda_u,z)\bigr).
\]
\end{definition}

\begin{proposition}[Evaluation is functorial on output-comparison cells]
\label{prop:relative-evaluation-functorial}
For every right \(\mathbb B\)-action \(Y\), the construction
\[
P\longmapsto \operatorname{Ev}_Y(P),
\qquad
\Gamma\longmapsto \operatorname{Ev}_Y(\Gamma)
\]
defines a functor
\[
\operatorname{Ev}_Y:\mathrm{Kl}(T_B)\to \mathcal E/A_0.
\]
\end{proposition}

\begin{proof}
\leavevmode
\begin{enumerate}
\item \textbf{Check well-definedness.}
If
\[
r(z)=q_Q(u)=t_B(\lambda_u),
\]
then \(\sigma(\lambda_u,z)\) is defined. Moreover
\[
r\sigma(\lambda_u,z)=s_B(\lambda_u)=q_P(\gamma u),
\]
so
\[
(\gamma u,\sigma(\lambda_u,z))
\]
lies in \(P\times_{B_0}Y\). It lies over \(A_0\) because
\[
p_P\gamma=p_Q.
\]

\item \textbf{Check identities.}
The Kleisli identity on \(P\) is
\[
x\longmapsto (x,e_B(q_Px)).
\]
Evaluating gives
\[
(x,z)\longmapsto
(x,\sigma(e_B(q_Px),z))
=
(x,z)
\]
by the unit law of the action \(Y\).

\item \textbf{Compute Kleisli composition.}
Let
\[
R\xrightarrow{\Delta}T_B(Q),
\qquad
Q\xrightarrow{\Gamma}T_B(P)
\]
be Kleisli arrows, with
\[
\Delta(w)=(\eta w,\mu_w),
\qquad
\Gamma(q)=(\gamma q,\lambda_q).
\]
The Kleisli composite has component
\[
w\longmapsto
\bigl(\gamma\eta w,\;m_B(\lambda_{\eta w},\mu_w)\bigr).
\]
In arrow notation this is the composite
\[
q_P(\gamma\eta w)
\xrightarrow{\lambda_{\eta w}}
q_Q(\eta w)
\xrightarrow{\mu_w}
q_R(w).
\]

\item \textbf{Evaluate the composite.}
Evaluation at \((w,z)\) gives
\[
\bigl(\gamma\eta w,
\sigma(m_B(\lambda_{\eta w},\mu_w),z)\bigr).
\]
By the right-action multiplication law, this equals
\[
\bigl(\gamma\eta w,
\sigma(\lambda_{\eta w},\sigma(\mu_w,z))\bigr),
\]
which is
\[
\operatorname{Ev}_Y(\Gamma)
\bigl(
\operatorname{Ev}_Y(\Delta)(w,z)
\bigr).
\]
\end{enumerate}
\end{proof}

\section{Restriction of downstream actions along Mealy morphisms}
\label{sec:relative-evaluation-mealy-morphisms}

Let
\[
\Phi=(P,\delta_P,\omega_P):\mathbb A\rightsquigarrow\mathbb B
\]
be a Mealy morphism. Thus \(P\) is a span
\[
\begin{tikzcd}
& P \arrow[dl,"p_P"'] \arrow[dr,"q_P"] & \\
A_0 && B_0
\end{tikzcd}
,
\]
and
\[
\delta_P:A_1\times_{t_A,A_0,p_P}P\to P,
\qquad
\omega_P:A_1\times_{t_A,A_0,p_P}P\to B_1
\]
satisfy
\[
p_P\delta_P=s_A\pi_1,
\qquad
s_B\omega_P=q_P\delta_P,
\qquad
t_B\omega_P=q_P\pi_P,
\]
together with the Mealy unit and multiplication laws.

Let
\[
Y=(Y\xrightarrow{r}B_0,\sigma)
\]
be a right \(\mathbb B\)-action.

\subsection{The restricted action}
\label{subsec:relative-restricted-action}

\begin{definition}[Restricted downstream action]
\label{def:relative-restricted-action}
The \textbf{restriction of \(Y\) along \(\Phi\)} is the right
\(\mathbb A\)-action
\[
\Phi^\ast Y
\]
whose underlying object over \(A_0\) is
\[
S_{\Phi,Y}:=
P\times_{q_P,B_0,r}Y
\xrightarrow{p_P\pi_P}
A_0,
\]
and whose action map is
\[
A_1\times_{t_A,A_0,p_P\pi_P}S_{\Phi,Y}
\to S_{\Phi,Y}
\]
given by
\[
a\cdot(x,y)
=
\bigl(
\delta_P(a,x),
\sigma(\omega_P(a,x),y)
\bigr).
\]
The action square has the form
\[
\begin{tikzcd}
A_1\times_{A_0}S_{\Phi,Y}
\arrow[rr,"{(a,x,y)\mapsto(\delta_P(a,x),\sigma(\omega_P(a,x),y))}"]
\arrow[dr,"s_A\pi_1"']
&&
S_{\Phi,Y}
\arrow[dl,"p_P\pi_P"]
\\
& A_0 &
\end{tikzcd}
.
\]
\end{definition}

\begin{theorem}[Mealy restriction of downstream actions]
\label{thm:relative-mealy-restriction-actions}
For every Mealy morphism
\[
\Phi:\mathbb A\rightsquigarrow\mathbb B
\]
and every right \(\mathbb B\)-action \(Y\), the object
\(\Phi^\ast Y\) of Definition~\ref{def:relative-restricted-action} is a right
\(\mathbb A\)-action.
\end{theorem}

\begin{proof}
\leavevmode
\begin{enumerate}
\item \textbf{Check that the formula lands in the pullback.}
If
\[
t_B\omega_P(a,x)=q_P(x)=r(y),
\]
then \(\sigma(\omega_P(a,x),y)\) is defined. Moreover
\[
r\sigma(\omega_P(a,x),y)
=
s_B\omega_P(a,x)
=
q_P\delta_P(a,x),
\]
so the displayed pair lies in
\[
P\times_{B_0}Y.
\]

\item \textbf{Check the projection to \(A_0\).}
The projection of the result to \(A_0\) is
\[
p_P\delta_P(a,x)=s_A(a).
\]
Thus the action has the required target-to-source typing.

\item \textbf{Check the unit law.}
Using the Mealy unit laws and the unit law of \(Y\),
\[
e_A(p_Px)\cdot(x,y)
=
\bigl(
\delta_P(e_A(p_Px),x),
\sigma(\omega_P(e_A(p_Px),x),y)
\bigr)
=
(x,y).
\]

\item \textbf{Check the multiplication law.}
Let
\[
a:u\to v,
\qquad
b:v\to w,
\qquad
p_P(x)=w.
\]
Then
\[
(b\circ a)\cdot(x,y)
=
\bigl(
\delta_P(b\circ a,x),
\sigma(\omega_P(b\circ a,x),y)
\bigr).
\]
The Mealy multiplication laws give
\[
\delta_P(b\circ a,x)
=
\delta_P(a,\delta_P(b,x))
\]
and
\[
\omega_P(b\circ a,x)
=
\omega_P(b,x)\circ\omega_P(a,\delta_P(b,x)).
\]
Equivalently,
\[
\omega_P(b\circ a,x)
=
m_B\bigl(
\omega_P(a,\delta_P(b,x)),
\omega_P(b,x)
\bigr).
\]
The right-action multiplication law for \(Y\) gives
\[
(b\circ a)\cdot(x,y)
=
a\cdot\bigl(b\cdot(x,y)\bigr).
\]
\end{enumerate}
\end{proof}

\subsection{Functoriality in the downstream action}
\label{subsec:relative-functoriality-downstream-action}

\begin{proposition}[Functoriality in the downstream action]
\label{prop:relative-functoriality-downstream-action}
For fixed \(\Phi:\mathbb A\rightsquigarrow\mathbb B\), the assignment
\[
Y\longmapsto \Phi^\ast Y
\]
extends to a functor
\[
\Phi^\ast:\mathbf{Act}(\mathbb B)\to\mathbf{Act}(\mathbb A).
\]
If
\[
f:Y\to Z
\]
is a \(\mathbb B\)-equivariant map, then
\[
\Phi^\ast f:S_{\Phi,Y}\to S_{\Phi,Z}
\]
is given by
\[
(x,y)\longmapsto(x,f(y)).
\]
\end{proposition}

\begin{proof}
\leavevmode
\begin{enumerate}
\item \textbf{Check well-definedness.}
The square
\[
\begin{tikzcd}
S_{\Phi,Y}
\arrow[r,"1_P\times f"]
\arrow[d]
&
S_{\Phi,Z}
\arrow[d]
\\
P
\arrow[r,"1_P"']
&
P
\end{tikzcd}
\]
is induced by the pullback defining \(S_{\Phi,Y}\), because \(f\) lies over
\(B_0\).

\item \textbf{Check the map over \(A_0\).}
The \(P\)-component is unchanged, so the map lies over \(A_0\).

\item \textbf{Check equivariance.}
For an admissible input arrow \(a\),
\[
f(\sigma_Y(\omega_P(a,x),y))
=
\sigma_Z(\omega_P(a,x),f(y))
\]
by equivariance of \(f\) as a map of right \(\mathbb B\)-actions.

\item \textbf{Check identity and composition.}
Both follow from the corresponding equations for maps in \(\mathcal E\).
\end{enumerate}
\end{proof}

\section{Evaluation of Mealy simulations}
\label{sec:relative-evaluation-mealy-simulations}

Let
\[
\Phi=(P,\delta_P,\omega_P),
\qquad
\Psi=(Q,\delta_Q,\omega_Q)
:
\mathbb A\rightsquigarrow\mathbb B
\]
be Mealy morphisms.

\subsection{Component form}
\label{subsec:relative-simulation-component-form}

\begin{definition}[Mealy simulation, component form]
\label{def:relative-mealy-simulation-component}
A \textbf{Mealy simulation}
\[
\Theta:\Phi\Rightarrow\Psi
\]
is a Kleisli output-comparison cell
\[
Q\to B\circ P.
\]
In components, write
\[
\Theta(y)=(\theta y,\alpha_y),
\]
where
\[
\theta:Q\to P,
\qquad
\alpha:Q\to B_1,
\]
and
\[
p_P\theta=p_Q,
\qquad
s_B\alpha=q_P\theta,
\qquad
t_B\alpha=q_Q.
\]
The typing data are summarized by
\[
\begin{tikzcd}[column sep=large,row sep=large]
&& B_1
  \arrow[dd,bend left=18,"s_B"]
  \arrow[dd,bend right=18,"t_B"']
&
\\
Q
  \arrow[rru,"\alpha"]
  \arrow[r,"\theta"']
  \arrow[drr,"q_Q"']
  \arrow[ddrr,bend right=16,"p_Q"']
&
P
  \arrow[dr,"q_P"]
  \arrow[rdd,bend left=16,"p_P"]
&&
\\
&& B_0 &
\\
&& A_0 &
\end{tikzcd}
.
\]
The simulation equations are
\[
\theta(\delta_Q(a,y))
=
\delta_P(a,\theta y),
\]
and
\[
\alpha_y\circ \omega_P(a,\theta y)
=
\omega_Q(a,y)\circ \alpha_{\delta_Q(a,y)}.
\]
Equivalently,
\[
m_B\bigl(\omega_P(a,\theta y),\alpha_y\bigr)
=
m_B\bigl(\alpha_{\delta_Q(a,y)},\omega_Q(a,y)\bigr).
\]
\end{definition}

\subsection{Evaluated simulations}
\label{subsec:relative-evaluated-simulations}

\begin{definition}[Evaluated simulation map]
\label{def:relative-evaluated-simulation-map}
Let \(Y\) be a right \(\mathbb B\)-action. The \textbf{evaluation} of
\(\Theta\) at \(Y\) is the map
\[
h_{\Theta,Y}:S_{\Psi,Y}\to S_{\Phi,Y}
\]
given by
\[
h_{\Theta,Y}(y,z)
=
\bigl(\theta y,\sigma(\alpha_y,z)\bigr).
\]
It is induced by the diagram
\[
\begin{tikzcd}
S_{\Psi,Y}
\arrow[r,"h_{\Theta,Y}"]
\arrow[d]
&
S_{\Phi,Y}
\arrow[d]
\\
Q
\arrow[r,"\theta"']
&
P
\end{tikzcd}
\]
together with the action of \(\alpha_y\) on the \(Y\)-coordinate.
\end{definition}

\begin{theorem}[Simulations evaluate to equivariant maps]
\label{thm:relative-simulations-evaluate-equivariantly}
For every Mealy simulation
\[
\Theta:\Phi\Rightarrow\Psi
\]
and every right \(\mathbb B\)-action \(Y\), the map
\[
h_{\Theta,Y}:S_{\Psi,Y}\to S_{\Phi,Y}
\]
is an \(\mathbb A\)-equivariant map
\[
\Psi^\ast Y\to \Phi^\ast Y.
\]
\end{theorem}

\begin{proof}
\leavevmode
\begin{enumerate}
\item \textbf{Check well-definedness.}
Since
\[
r\sigma(\alpha_y,z)=s_B\alpha_y=q_P(\theta y),
\]
the pair
\[
(\theta y,\sigma(\alpha_y,z))
\]
lies in \(S_{\Phi,Y}\).

\item \textbf{Check the map over \(A_0\).}
The equation
\[
p_P\theta=p_Q
\]
gives
\[
p_P(\theta y)=p_Q(y).
\]

\item \textbf{Compute the left-hand equivariance composite.}
Let
\[
a:u\to v,
\qquad
p_Q(y)=v,
\qquad
r(z)=q_Q(y).
\]
Then
\[
h_{\Theta,Y}\bigl(a\cdot_\Psi(y,z)\bigr)
=
h_{\Theta,Y}
\bigl(
\delta_Q(a,y),
\sigma(\omega_Q(a,y),z)
\bigr),
\]
which is
\[
\bigl(
\theta\delta_Q(a,y),
\sigma(\alpha_{\delta_Q(a,y)},\sigma(\omega_Q(a,y),z))
\bigr).
\]
Using the action multiplication law for \(Y\), the second component is
\[
\sigma\bigl(
\omega_Q(a,y)\circ \alpha_{\delta_Q(a,y)},
z
\bigr).
\]

\item \textbf{Use the simulation law.}
The simulation law gives
\[
\omega_Q(a,y)\circ \alpha_{\delta_Q(a,y)}
=
\alpha_y\circ \omega_P(a,\theta y).
\]
Thus the second component becomes
\[
\sigma\bigl(
\alpha_y\circ \omega_P(a,\theta y),
z
\bigr).
\]
Using the action multiplication law again, this is
\[
\sigma(\omega_P(a,\theta y),\sigma(\alpha_y,z)).
\]

\item \textbf{Compare with the right-hand equivariance composite.}
The state equation
\[
\theta\delta_Q(a,y)=\delta_P(a,\theta y)
\]
gives
\[
h_{\Theta,Y}\bigl(a\cdot_\Psi(y,z)\bigr)
=
a\cdot_\Phi h_{\Theta,Y}(y,z).
\]
\end{enumerate}
\end{proof}

\begin{corollary}[Contravariance in simulations]
\label{cor:relative-contravariance-actions}
For fixed \(Y\in\mathbf{Act}(\mathbb B)\), the assignment
\[
\Phi\longmapsto\Phi^\ast Y
\]
is contravariant in Mealy simulations:
\[
\Theta:\Phi\Rightarrow\Psi
\quad\longmapsto\quad
h_{\Theta,Y}:\Psi^\ast Y\to\Phi^\ast Y.
\]
Thus it defines a functor
\[
\mathbf{Mly}(\mathcal E)(\mathbb A,\mathbb B)^{\mathrm{op}}
\longrightarrow
\mathbf{Act}(\mathbb A).
\]
Here \(\mathbf{Mly}(\mathcal E)(\mathbb A,\mathbb B)\) is the hom-category
with the displayed simulation direction of
Chapter~\ref{ch:spans-internal-categories-mealy}.
\end{corollary}

\begin{proof}
\leavevmode
\begin{enumerate}
\item \textbf{Check identities.}
The identity simulation uses the identity output arrows of \(\mathbb B\).
Evaluation sends these identity output arrows to identity transitions in \(Y\).

\item \textbf{Check composition.}
Composition of simulations is Kleisli composition. By
Proposition~\ref{prop:relative-evaluation-functorial}, downstream evaluation
preserves Kleisli composition.

\item \textbf{Check the direction.}
A simulation
\[
\Theta:\Phi\Rightarrow\Psi
\]
is represented by a Kleisli cell from the carrier of \(\Psi\) to the
output-comparison object of the carrier of \(\Phi\). Hence evaluation produces
a map
\[
\Psi^\ast Y\to\Phi^\ast Y.
\]
\end{enumerate}
\end{proof}

\section{Relative Mealy program categories}
\label{sec:relative-program-categories}

Fix a Mealy morphism
\[
\Phi=(P,\delta_P,\omega_P):\mathbb A\rightsquigarrow\mathbb B
\]
and a right \(\mathbb B\)-action
\[
Y=(Y\xrightarrow{r}B_0,\sigma).
\]

\subsection{State and primitive execution objects}
\label{subsec:relative-state-primitive-objects}

\begin{definition}[Relative state and primitive execution objects]
\label{def:relative-state-and-primitive}
The \textbf{relative state object} is
\[
S_{\Phi,Y}:=P\times_{q_P,B_0,r}Y.
\]
A generalized element is written
\[
(x,y),
\qquad
q_P(x)=r(y).
\]
The \textbf{relative primitive execution object} is
\[
M_{\Phi,Y}
:=
A_1\times_{t_A,A_0,p_P\pi_P}S_{\Phi,Y}.
\]
It is defined by the pullback
\[
\begin{tikzcd}
M_{\Phi,Y}
\arrow[r,"\pi_S"]
\arrow[d,"\pi_A"']
&
S_{\Phi,Y}
\arrow[d,"p_P\pi_P"]
\\
A_1
\arrow[r,"t_A"']
&
A_0
\arrow[ul,phantom,very near start,"\lrcorner"]
\end{tikzcd}
.
\]
A generalized element is written
\[
(a,x,y),
\qquad
t_A(a)=p_P(x),
\qquad
q_P(x)=r(y).
\]
\end{definition}

\subsection{The program category}
\label{subsec:relative-program-category-definition}

\begin{definition}[Relative Mealy program category]
\label{def:relative-mealy-program-category}
The \textbf{relative Mealy program category} of \(\Phi\) at \(Y\) is
\[
\mathsf{Prog}_{\Phi,Y}
:=
\int_{\mathbb A}\Phi^\ast Y.
\]
Explicitly,
\[
(\mathsf{Prog}_{\Phi,Y})_0=S_{\Phi,Y},
\qquad
(\mathsf{Prog}_{\Phi,Y})_1=M_{\Phi,Y}.
\]
The source and target maps are
\[
s_{\Phi,Y}(a,x,y)=(x,y),
\]
and
\[
t_{\Phi,Y}(a,x,y)
=
\bigl(
\delta_P(a,x),
\sigma(\omega_P(a,x),y)
\bigr).
\]
They form the primitive execution span
\[
\mathbf m_{\Phi,Y}
=
\left(
S_{\Phi,Y}
\xleftarrow{s_{\Phi,Y}}
M_{\Phi,Y}
\xrightarrow{t_{\Phi,Y}}
S_{\Phi,Y}
\right).
\]
The identity is
\[
e_{\Phi,Y}(x,y)=(e_A(p_Px),x,y).
\]
Composition is
\[
\bigl(
(b,x,y),
(a,\delta_P(b,x),\sigma(\omega_P(b,x),y))
\bigr)
\longmapsto
(m_A(a,b),x,y).
\]
\end{definition}

\begin{theorem}[Relative program categories]
\label{thm:relative-program-categories}
For every Mealy morphism
\[
\Phi:\mathbb A\rightsquigarrow\mathbb B
\]
and every right \(\mathbb B\)-action \(Y\), the data of
Definition~\ref{def:relative-mealy-program-category} define an internal
category
\[
\mathsf{Prog}_{\Phi,Y}.
\]
Equivalently, the span
\[
\mathbf m_{\Phi,Y}
=
\left(
S_{\Phi,Y}
\xleftarrow{s_{\Phi,Y}}
M_{\Phi,Y}
\xrightarrow{t_{\Phi,Y}}
S_{\Phi,Y}
\right)
\]
is a monoid object in
\[
\operatorname{EndSpan}_{\mathcal E}(S_{\Phi,Y}).
\]
\end{theorem}

\begin{proof}
\leavevmode
\begin{enumerate}
\item \textbf{Use the restricted action.}
By Theorem~\ref{thm:relative-mealy-restriction-actions},
\(\Phi^\ast Y\) is a right \(\mathbb A\)-action.

\item \textbf{Apply the action-category construction.}
Theorem~\ref{thm:relative-action-category-internal} applied to this action
gives the internal category
\[
\int_{\mathbb A}\Phi^\ast Y.
\]

\item \textbf{Identify the displayed formulas.}
The source, target, identity, and composition maps in
Definition~\ref{def:relative-mealy-program-category} are precisely the
specialized action-category maps.
\end{enumerate}
\end{proof}

\subsection{Program functors induced by simulations}
\label{subsec:relative-program-functors-simulations}

\begin{definition}[Program functor induced by a simulation]
\label{def:relative-program-functor-simulation}
Let
\[
\Theta:\Phi\Rightarrow\Psi
\]
be a Mealy simulation with components
\[
\theta:Q\to P,
\qquad
\alpha:Q\to B_1.
\]
For every right \(\mathbb B\)-action \(Y\), define
\[
H_{\Theta,Y}:
\mathsf{Prog}_{\Psi,Y}
\to
\mathsf{Prog}_{\Phi,Y}
\]
on objects by
\[
H_{\Theta,Y}(y,z)
=
\bigl(\theta y,\sigma(\alpha_y,z)\bigr),
\]
and on arrows by
\[
H_{\Theta,Y}(a,y,z)
=
\bigl(a,\theta y,\sigma(\alpha_y,z)\bigr).
\]
The arrow part is induced by the pullback square
\[
\begin{tikzcd}
M_{\Psi,Y}
\arrow[r,"k_{\Theta,Y}"]
\arrow[d,"\pi_S"']
&
M_{\Phi,Y}
\arrow[d,"\pi_S"]
\\
S_{\Psi,Y}
\arrow[r,"h_{\Theta,Y}"']
&
S_{\Phi,Y}
\end{tikzcd}
\]
because both primitive execution objects are pulled back from \(A_1\).
\end{definition}

\begin{theorem}[Simulations induce contravariant program functors]
\label{thm:relative-simulations-program-functors}
The maps in Definition~\ref{def:relative-program-functor-simulation} define an
internal functor
\[
H_{\Theta,Y}:
\mathsf{Prog}_{\Psi,Y}
\to
\mathsf{Prog}_{\Phi,Y}.
\]
Consequently, for fixed \(Y\), the assignment
\[
\Phi\longmapsto\mathsf{Prog}_{\Phi,Y}
\]
extends to a functor
\[
\mathbf{Mly}(\mathcal E)(\mathbb A,\mathbb B)^{\mathrm{op}}
\longrightarrow
\mathbf{Cat}_{\mathrm{int}}(\mathcal E)/\mathbb A^{\mathrm{op}}.
\]
\end{theorem}

\begin{proof}
\leavevmode
\begin{enumerate}
\item \textbf{Identify the object map.}
The object map is
\[
h_{\Theta,Y}:S_{\Psi,Y}\to S_{\Phi,Y}.
\]

\item \textbf{Identify the arrow map.}
The arrow map is the induced map on the pullbacks defining primitive execution
objects:
\[
(a,y,z)\longmapsto
(a,\theta y,\sigma(\alpha_y,z)).
\]

\item \textbf{Check source and target preservation.}
Source preservation is immediate. Target preservation is exactly the
equivariance of \(h_{\Theta,Y}\), proved in
Theorem~\ref{thm:relative-simulations-evaluate-equivariantly}.

\item \textbf{Check identities and composition.}
These follow from functoriality of action categories with respect to
equivariant maps.

\item \textbf{Check the slice condition.}
The functor lands over \(\mathbb A^{\mathrm{op}}\) because the arrow map leaves
the \(A_1\)-component unchanged and the object map lies over \(A_0\).
\end{enumerate}
\end{proof}

\section{Input and downstream comparison functors}
\label{sec:relative-input-downstream-comparison}

Let
\[
\Phi=(P,\delta_P,\omega_P):\mathbb A\rightsquigarrow\mathbb B
\]
be a Mealy morphism, and let
\[
Y=(Y\xrightarrow{r}B_0,\sigma)
\]
be a right \(\mathbb B\)-action.

\subsection{Input projection}
\label{subsec:relative-input-projection}

\begin{definition}[Input projection]
\label{def:relative-input-projection}
The \textbf{input projection}
\[
I_{\Phi,Y}:\mathsf{Prog}_{\Phi,Y}\to\mathbb A^{\mathrm{op}}
\]
is defined on objects by
\[
I_{\Phi,Y}(x,y)=p_P(x),
\]
and on arrows by
\[
I_{\Phi,Y}(a,x,y)=a.
\]
It is displayed by
\[
\begin{tikzcd}
M_{\Phi,Y}
\arrow[r,"\pi_A"]
\arrow[d,"s_{\Phi,Y}"']
&
A_1
\arrow[d,"t_A"]
\\
S_{\Phi,Y}
\arrow[r,"p_P\pi_P"']
&
A_0
\end{tikzcd}
.
\]
\end{definition}

\subsection{Downstream comparison}
\label{subsec:relative-downstream-comparison}

\begin{definition}[Downstream comparison functor]
\label{def:relative-downstream-comparison}
The \textbf{downstream comparison functor}
\[
\Lambda_{\Phi,Y}:\mathsf{Prog}_{\Phi,Y}\to\int_{\mathbb B}Y
\]
is defined on objects by
\[
\Lambda_{\Phi,Y}(x,y)=y,
\]
and on arrows by
\[
\Lambda_{\Phi,Y}(a,x,y)
=
(\omega_P(a,x),y).
\]
The arrow map is typed by the pullback
\[
\begin{tikzcd}
M_{\Phi,Y}
\arrow[rr,"{(a,x,y)\mapsto(\omega_P(a,x),y)}"]
\arrow[dr,"\pi_Y\pi_S"']
&&
B_1\times_{t_B,B_0,r}Y
\arrow[dl,"\pi_Y"]
\\
& Y &
\end{tikzcd}
.
\]
\end{definition}

\begin{theorem}[Input and downstream functors]
\label{thm:relative-input-downstream-functors}
The input projection
\[
I_{\Phi,Y}:\mathsf{Prog}_{\Phi,Y}\to\mathbb A^{\mathrm{op}}
\]
is input-discrete. The downstream comparison
\[
\Lambda_{\Phi,Y}:\mathsf{Prog}_{\Phi,Y}\to\int_{\mathbb B}Y
\]
is an internal functor.
\end{theorem}

\begin{proof}
\leavevmode
\begin{enumerate}
\item \textbf{Check input-discreteness.}
The input projection is the canonical projection associated to the right
\(\mathbb A\)-action \(\Phi^\ast Y\). Hence it is input-discrete by
Theorem~\ref{thm:relative-actions-input-discrete}.

\item \textbf{Check that the input projection is a functor.}
The source square is the defining pullback. The target square commutes because
\[
p_P\delta_P(a,x)=s_A(a),
\]
which is the target equation for a functor into \(\mathbb A^{\mathrm{op}}\).
Identity and composition preservation follow from the identity and
multiplication laws of \(\mathbb A\).

\item \textbf{Check the downstream arrow map.}
The arrow map of \(\Lambda_{\Phi,Y}\) is well-defined because
\[
t_B\omega_P(a,x)=q_P(x)=r(y).
\]

\item \textbf{Check source and target preservation.}
In \(\int_{\mathbb B}Y\),
\[
s(\omega_P(a,x),y)=y,
\]
and
\[
t(\omega_P(a,x),y)=\sigma(\omega_P(a,x),y).
\]
These are the source and the \(Y\)-component of the target in
\(\mathsf{Prog}_{\Phi,Y}\).

\item \textbf{Check identities and composition.}
Identity preservation follows from the Mealy unit law for \(\omega_P\). For
composition, let
\[
a:u\to v,
\qquad
b:v\to w,
\qquad
p_P(x)=w.
\]
The composite of the two downstream arrows has \(B\)-label
\[
m_B\bigl(\omega_P(a,\delta_P(b,x)),\omega_P(b,x)\bigr)
=
\omega_P(b,x)\circ \omega_P(a,\delta_P(b,x)).
\]
By the Mealy multiplication law, this equals
\[
\omega_P(b\circ a,x).
\]
Thus \(\Lambda_{\Phi,Y}\) preserves composition.
\end{enumerate}
\end{proof}

\begin{corollary}[Combined labelling]
\label{cor:relative-combined-labelling}
Since \(\mathcal E\) has finite limits, the product internal category
\[
\mathbb A^{\mathrm{op}}\times\int_{\mathbb B}Y
\]
exists. Hence \(I_{\Phi,Y}\) and \(\Lambda_{\Phi,Y}\) assemble into an internal
functor
\[
L_{\Phi,Y}:
\mathsf{Prog}_{\Phi,Y}
\to
\mathbb A^{\mathrm{op}}\times\int_{\mathbb B}Y.
\]
On objects,
\[
L_{\Phi,Y}(x,y)=(p_Px,y),
\]
and on arrows,
\[
L_{\Phi,Y}(a,x,y)=\bigl(a,(\omega_P(a,x),y)\bigr).
\]
\end{corollary}

\begin{proof}
\leavevmode
\begin{enumerate}
\item \textbf{Use the product universal property.}
The functors \(I_{\Phi,Y}\) and \(\Lambda_{\Phi,Y}\) have common domain
\(\mathsf{Prog}_{\Phi,Y}\). The product internal category therefore supplies
a unique induced functor.

\item \textbf{Identify the components.}
The displayed object and arrow maps are the two projections of this induced
functor.
\end{enumerate}
\end{proof}

\section{Output-comparison transformations after evaluation}
\label{sec:relative-output-comparison-transformations}

Let
\[
\Theta:\Phi\Rightarrow\Psi
\]
be a Mealy simulation with components
\[
\theta:Q\to P,
\qquad
\alpha:Q\to B_1.
\]
Let \(Y\) be a right \(\mathbb B\)-action.

The induced functor
\[
H_{\Theta,Y}:
\mathsf{Prog}_{\Psi,Y}\to\mathsf{Prog}_{\Phi,Y}
\]
strictly preserves input labels:
\[
I_{\Phi,Y}H_{\Theta,Y}=I_{\Psi,Y}.
\]
The output-comparison component \(\alpha\) becomes a natural transformation in
the downstream action category.

\subsection{The evaluated downstream transformation}
\label{subsec:relative-downstream-transformation}

\begin{definition}[Evaluated downstream comparison transformation]
\label{def:relative-downstream-comparison-transformation}
Define
\[
\lambda_{\Theta,Y}:
\Lambda_{\Psi,Y}
\Rightarrow
\Lambda_{\Phi,Y}H_{\Theta,Y}
\]
by assigning to a state
\[
(y,z)\in S_{\Psi,Y}
\]
the arrow
\[
(\alpha_y,z)
:
z\to \sigma(\alpha_y,z)
\]
of the action category \(\int_{\mathbb B}Y\).
The component lies in
\[
B_1\times_{t_B,B_0,r}Y
\]
because
\[
t_B\alpha_y=q_Q(y)=r(z).
\]
Its source and target are
\[
\pi_Y(\alpha_y,z)=z,
\qquad
\sigma(\alpha_y,z).
\]
The relevant typing equations for the action-arrow object are
\[
r\pi_Y=t_B\pi_{B_1},
\qquad
r\sigma=s_B\pi_{B_1}.
\]
Equivalently, the component is represented by the diagram
\[
\begin{tikzcd}[column sep=large,row sep=large]
&
B_1\times_{t_B,B_0,r}Y
\arrow[dl,"\pi_Y"']
\arrow[dr,"\sigma"]
\arrow[d,"\pi_{B_1}" description]
&
\\
Y
\arrow[dr,"r"']
&
B_1
\arrow[d,"{(t_B,s_B)}" description]
&
Y
\arrow[dl,"r"]
\\
&
B_0\times B_0
&
\end{tikzcd}
.
\]
Thus \((\alpha_y,z)\) is viewed as a right-action arrow from the object over
\(q_Q(y)\) to the object over \(q_P(\theta y)\).
\end{definition}

\begin{theorem}[Naturality of the evaluated output comparison]
\label{thm:relative-output-comparison-natural}
The family
\[
\lambda_{\Theta,Y}
\]
of Definition~\ref{def:relative-downstream-comparison-transformation} is an
internal natural transformation
\[
\Lambda_{\Psi,Y}
\Rightarrow
\Lambda_{\Phi,Y}H_{\Theta,Y}.
\]
\end{theorem}

\begin{proof}
\leavevmode
\begin{enumerate}
\item \textbf{Check the components.}
The component at \((y,z)\) is well-defined because
\[
t_B\alpha_y=q_Q(y)=r(z).
\]
Its source is \(z\), and its target is
\[
\sigma(\alpha_y,z),
\]
which is the object part of
\[
\Lambda_{\Phi,Y}H_{\Theta,Y}(y,z).
\]

\item \textbf{Write the naturality square.}
Let
\[
(a,y,z):(y,z)\to
\bigl(\delta_Q(a,y),\sigma(\omega_Q(a,y),z)\bigr)
\]
be a primitive arrow of \(\mathsf{Prog}_{\Psi,Y}\). Naturality in
\(\int_{\mathbb B}Y\) asks for the commutativity of
\[
\begin{tikzcd}[column sep=large,row sep=large]
z
  \arrow[r,"{(\omega_Q(a,y),z)}"]
  \arrow[d,"{(\alpha_y,z)}"']
&
\sigma(\omega_Q(a,y),z)
  \arrow[d,"{(\alpha_{\delta_Q(a,y)},\,\sigma(\omega_Q(a,y),z))}"]
\\
\sigma(\alpha_y,z)
  \arrow[r,"{(\omega_P(a,\theta y),\,\sigma(\alpha_y,z))}"']
&
\sigma(c_{\Theta,a,y},z)
\end{tikzcd}
\]
where
\[
c_{\Theta,a,y}
:=
\omega_Q(a,y)\circ \alpha_{\delta_Q(a,y)}
=
\alpha_y\circ \omega_P(a,\theta y).
\]

\item \textbf{Make the composition order explicit.}
With the convention
\[
m_B(\beta,\gamma)=\gamma\circ\beta,
\]
the top-right composite corresponds to
\[
m_B\bigl(\alpha_{\delta_Q(a,y)},\omega_Q(a,y)\bigr)
=
\omega_Q(a,y)\circ\alpha_{\delta_Q(a,y)}.
\]
The left-bottom composite corresponds to
\[
m_B\bigl(\omega_P(a,\theta y),\alpha_y\bigr)
=
\alpha_y\circ\omega_P(a,\theta y).
\]

\item \textbf{Use the simulation law.}
The equality of these two composites is exactly the output equation in the
Mealy simulation law:
\[
\alpha_y\circ \omega_P(a,\theta y)
=
\omega_Q(a,y)\circ \alpha_{\delta_Q(a,y)}.
\]
\end{enumerate}
\end{proof}

\begin{corollary}[Relative labelled functor induced by a simulation]
\label{cor:relative-labelled-functor-simulation}
Every Mealy simulation
\[
\Theta:\Phi\Rightarrow\Psi
\]
and every right \(\mathbb B\)-action \(Y\) induce an input-strict,
downstream-lax labelled functor
\[
H_{\Theta,Y}:
(\mathsf{Prog}_{\Psi,Y},I_{\Psi,Y},\Lambda_{\Psi,Y})
\longrightarrow
(\mathsf{Prog}_{\Phi,Y},I_{\Phi,Y},\Lambda_{\Phi,Y})
\]
consisting of
\[
I_{\Phi,Y}H_{\Theta,Y}=I_{\Psi,Y}
\]
and a canonical natural transformation
\[
\lambda_{\Theta,Y}:
\Lambda_{\Psi,Y}
\Rightarrow
\Lambda_{\Phi,Y}H_{\Theta,Y}.
\]
\end{corollary}

\section{Predicate transformers for spans}
\label{sec:relative-predicate-transformers}

We now pass from the operational layer to the logical layer. The notation is
chosen to match the non-iterative dynamic-logic reading of programs as
predicate transformers: diamonds express existence of a successful execution,
boxes express universal correctness over all executions, tests restrict
execution to states satisfying a predicate, and finite choice denotes
nondeterministic branching \cite{HarelKozenTiurynDynamicLogic}. The
categorical implementation uses subobjects as predicates, regular images as
existential quantification, and right adjoints to pullback as universal
quantification \cite{Jacobs1999CategoricalLogic}. In the regular fragment,
this agrees with the usual relational reading of programs, where relations can
be represented categorically by suitable spans
\cite{FongSpivakRegularRelational}.

\subsection{Regular images and existential quantification}
\label{subsec:relative-regular-images}

\begin{definition}[Regular category]
\label{def:relative-regular-category}
A category \(\mathcal E\) is \textbf{regular} if it has finite limits, every
morphism factors as a regular epimorphism followed by a monomorphism, and
regular epimorphisms are stable under pullback.
\end{definition}

For the rest of this section, assume \(\mathcal E\) is regular. For an object
\(X\), write
\[
\mathrm{Sub}(X)
\]
for the poset of subobjects of \(X\). For a morphism \(f:X\to Y\), write
\[
f^\ast:\mathrm{Sub}(Y)\to\mathrm{Sub}(X)
\]
for pullback.

\begin{definition}[Existential image]
\label{def:relative-existential-image}
For a morphism
\[
f:X\to Y,
\]
define
\[
\exists_f:\mathrm{Sub}(X)\to\mathrm{Sub}(Y)
\]
by sending a subobject \(U\hookrightarrow X\) to the regular image of
\[
U\to X\xrightarrow{f}Y.
\]
\end{definition}

\begin{proposition}[Existential adjunction]
\label{prop:relative-existential-adjunction}
For every morphism \(f:X\to Y\) in a regular category,
\[
\exists_f\dashv f^\ast.
\]
Thus, for \(U\in\mathrm{Sub}(X)\) and \(V\in\mathrm{Sub}(Y)\),
\[
\exists_f(U)\le V
\quad\Longleftrightarrow\quad
U\le f^\ast V.
\]
\end{proposition}

\begin{proof}
\leavevmode
\begin{enumerate}
\item \textbf{Translate the left-hand inequality.}
The statement
\[
\exists_f(U)\le V
\]
says that the image of the composite
\[
U\to X\xrightarrow{f}Y
\]
factors through \(V\hookrightarrow Y\).

\item \textbf{Remove the image.}
This is equivalent to the composite \(U\to Y\) itself factoring through
\(V\hookrightarrow Y\).

\item \textbf{Use the pullback.}
Factoring through \(V\) is equivalent to \(U\to X\) factoring through
\(f^\ast V\hookrightarrow X\).
\end{enumerate}
\end{proof}

\begin{proposition}[Beck--Chevalley for regular images]
\label{prop:relative-regular-beck-chevalley}
For every pullback square
\[
\begin{tikzcd}
X' \arrow[r,"g'"] \arrow[d,"f'"']
&
X \arrow[d,"f"]
\\
Y' \arrow[r,"g"']
&
Y
\arrow[ul,phantom,very near start,"\lrcorner"]
\end{tikzcd}
\]
in a regular category,
\[
g^\ast\exists_f=\exists_{f'}(g')^\ast.
\]
\end{proposition}

\begin{proof}
\leavevmode
\begin{enumerate}
\item \textbf{Factor the relevant composite.}
The regular image of a subobject after applying \(f\) is represented by a
regular epimorphism followed by a monomorphism.

\item \textbf{Pull the factorization back.}
Regular epimorphisms are stable under pullback, so pulling this image
factorization back along \(g\) gives the image factorization of the pulled
back subobject along \(f'\).

\item \textbf{Compare the two subobjects.}
The two constructions therefore determine the same subobject of \(Y'\).
\end{enumerate}
\end{proof}

When
\[
R:X\nrightarrow Y
\qquad\text{and}\qquad
S:Y\nrightarrow Z
\]
are spans, we write
\[
R;S:X\nrightarrow Z
\]
for the composite span obtained by performing \(R\) first and \(S\) second.
Equivalently, in the notation of
Chapter~\ref{ch:spans-internal-categories-mealy},
\[
R;S=S\circ R.
\]
This convention is the dynamic-logic convention for sequential composition of
programs.

\subsection{Diamond modalities}
\label{subsec:relative-diamond-modalities}

\begin{definition}[Diamond modality of a span]
\label{def:relative-span-diamond}
Let
\[
R=
\left(
X\xleftarrow{\ell}R\xrightarrow{r}Y
\right)
\]
be a span in a regular category. Define
\[
\langle R\rangle:\mathrm{Sub}(Y)\to\mathrm{Sub}(X)
\]
by
\[
\langle R\rangle\varphi
:=
\exists_\ell(r^\ast\varphi).
\]
The construction is represented by
\[
\begin{tikzcd}
r^\ast\varphi
\arrow[r]
\arrow[d]
&
\varphi
\arrow[d,hook]
\\
R
\arrow[r,"r"']
\arrow[d,"\ell"]
&
Y
\\
X
\end{tikzcd}
.
\]
\end{definition}

\begin{proposition}[Functoriality of diamonds]
\label{prop:relative-diamond-functoriality}
For spans
\[
R:X\nrightarrow Y
\qquad\text{and}\qquad
S:Y\nrightarrow Z
\]
in a regular category,
\[
\langle 1_X\rangle=\mathrm{id}_{\mathrm{Sub}(X)}
\]
and
\[
\langle R;S\rangle
=
\langle R\rangle\circ\langle S\rangle.
\]
\end{proposition}

\begin{proof}
\leavevmode
\begin{enumerate}
\item \textbf{Check the identity span.}
For the identity span
\[
X\xleftarrow{1_X}X\xrightarrow{1_X}X,
\]
the formula gives
\[
\exists_{1_X}1_X^\ast=\mathrm{id}_{\mathrm{Sub}(X)}.
\]

\item \textbf{Display the composite span.}
Write
\[
R=(X\xleftarrow{\ell}R\xrightarrow{r}Y),
\qquad
S=(Y\xleftarrow{m}S\xrightarrow{n}Z).
\]
Their composite has apex \(R\times_Y S\):
\[
\begin{tikzcd}[column sep=large,row sep=large]
R\times_Y S
  \arrow[r,"\pi_S"]
  \arrow[d,"\pi_R"']
&
S
  \arrow[d,"m"]
  \arrow[dr,"n",bend left=12]
&
\\
R
  \arrow[r,"r"']
  \arrow[d,"\ell"']
&
Y
&
Z
\\
X
&&
\end{tikzcd}
.
\]

\item \textbf{Compute the diamond of the composite.}
For \(\varphi\in\mathrm{Sub}(Z)\),
\[
\langle R;S\rangle\varphi
=
\exists_{\ell\pi_R}\bigl((n\pi_S)^\ast\varphi\bigr).
\]

\item \textbf{Use Beck--Chevalley.}
By functoriality of existential images and Beck--Chevalley for the pullback
defining \(R\times_Y S\), this is
\[
\exists_\ell r^\ast\exists_m n^\ast\varphi
=
\langle R\rangle(\langle S\rangle\varphi).
\]
\end{enumerate}
\end{proof}

\section{Relative primitive modalities}
\label{sec:relative-primitive-modalities}

Fix
\[
\Phi:\mathbb A\rightsquigarrow\mathbb B
\]
and a right \(\mathbb B\)-action \(Y\). Write
\[
S:=S_{\Phi,Y},
\qquad
M:=M_{\Phi,Y},
\]
and
\[
\mathbf m_{\Phi,Y}
=
(S\xleftarrow{s}M\xrightarrow{t}S)
\]
for the primitive execution span.

\subsection{Primitive diamonds}
\label{subsec:relative-primitive-diamonds}

\begin{definition}[Primitive diamond]
\label{def:relative-primitive-diamond}
For a predicate
\[
\varphi\hookrightarrow S,
\]
the \textbf{primitive diamond} is
\[
\langle \mathbf m_{\Phi,Y}\rangle\varphi
=
\exists_s(t^\ast\varphi).
\]
\end{definition}

\begin{proposition}[External reading of primitive diamonds]
\label{prop:relative-primitive-diamond-set-reading}
In \(\mathbf{Set}\),
\[
(x,y)\models \langle \mathbf m_{\Phi,Y}\rangle\varphi
\]
if and only if there exists an admissible input arrow
\[
a:u\to p_P(x)
\]
such that
\[
\bigl(
\delta_P(a,x),
\sigma(\omega_P(a,x),y)
\bigr)
\models \varphi.
\]
\end{proposition}

\begin{proof}
\leavevmode
\begin{enumerate}
\item \textbf{Expand the existential image.}
A witness for
\[
(x,y)\in\exists_s(t^\ast\varphi)
\]
is an element of \(M_{\Phi,Y}\) with source \((x,y)\).

\item \textbf{Identify such witnesses.}
Such an element is exactly a primitive execution
\[
(a,x,y)
\]
where \(a:u\to p_P(x)\).

\item \textbf{Read the target condition.}
The condition that this witness lies in \(t^\ast\varphi\) says that
\[
\bigl(
\delta_P(a,x),
\sigma(\omega_P(a,x),y)
\bigr)
\models\varphi.
\]
\end{enumerate}
\end{proof}

\subsection{Tests}
\label{subsec:relative-tests}

\begin{definition}[Tests]
\label{def:relative-tests}
For a state predicate
\[
\theta:\Theta\hookrightarrow S,
\]
define the corresponding test span
\[
\theta?
=
(S\xleftarrow{\theta}\Theta\xrightarrow{\theta}S).
\]
It is displayed as
\[
\begin{tikzcd}
& \Theta \arrow[dl,"\theta"'] \arrow[dr,"\theta"] & \\
S && S
\end{tikzcd}
.
\]
\end{definition}

\begin{proposition}[Diamond law for tests]
\label{prop:relative-test-diamond}
For every predicate \(\varphi\hookrightarrow S\),
\[
\langle\theta?\rangle\varphi
=
\theta\wedge\varphi.
\]
\end{proposition}

\begin{proof}
\leavevmode
\begin{enumerate}
\item \textbf{Expand the diamond.}
By definition,
\[
\langle\theta?\rangle\varphi
=
\exists_\theta(\theta^\ast\varphi).
\]

\item \textbf{Use monicity of the test.}
Since \(\theta\) is monic, existential image along \(\theta\) identifies
subobjects of \(\Theta\) with subobjects of \(S\) lying below \(\theta\).

\item \textbf{Identify the image.}
The image of \(\theta^\ast\varphi\) is therefore the meet
\[
\theta\wedge\varphi.
\]
\end{enumerate}
\end{proof}

\section{Restricted primitive subprograms}
\label{sec:relative-restricted-subprograms}

Let
\[
\iota_A:M_{\Phi,Y}\to A_1
\]
be the input projection
\[
(a,x,y)\mapsto a.
\]
Let
\[
\iota_B:M_{\Phi,Y}\to B_1
\]
be the emitted-output projection
\[
(a,x,y)\mapsto \omega_P(a,x).
\]
Let
\[
\Lambda_{\Phi,Y,1}:M_{\Phi,Y}\to
\left(\int_{\mathbb B}Y\right)_1
\]
be the arrow map of the downstream comparison functor:
\[
(a,x,y)\mapsto(\omega_P(a,x),y).
\]
These maps are related by
\[
\begin{tikzcd}[column sep=large]
M_{\Phi,Y}
\arrow[r,"\Lambda_{\Phi,Y,1}"]
\arrow[dr,"\iota_B"']
&
\left(\int_{\mathbb B}Y\right)_1
=
B_1\times_{t_B,B_0,r}Y
\arrow[d,"\pi_{B_1}"]
\\
&
B_1
\end{tikzcd}
.
\]

\begin{definition}[Restricted primitive subprogram]
\label{def:relative-restricted-primitive-subprogram}
Let
\[
R\hookrightarrow A_1
\]
be an input predicate, let
\[
O\hookrightarrow B_1
\]
be an emitted-output predicate, and let
\[
T\hookrightarrow \left(\int_{\mathbb B}Y\right)_1
\]
be a downstream-transition predicate. Define
\[
R^\sharp:=\iota_A^\ast R,
\qquad
O^\sharp:=\iota_B^\ast O,
\qquad
T^\sharp:=\Lambda_{\Phi,Y,1}^\ast T.
\]
Their simultaneous restriction is
\[
(R/O/T)^\sharp
:=
R^\sharp\wedge O^\sharp\wedge T^\sharp
\hookrightarrow M_{\Phi,Y}.
\]
The associated restricted primitive span is
\[
S_{\Phi,Y}
\xleftarrow{s i}
(R/O/T)^\sharp
\xrightarrow{t i}
S_{\Phi,Y},
\]
where
\[
i:(R/O/T)^\sharp\hookrightarrow M_{\Phi,Y}
\]
is the inclusion. Its diamond is denoted
\[
\langle R/O/T\rangle\varphi.
\]
\end{definition}

\begin{proposition}[External reading of restricted primitive diamonds]
\label{prop:relative-restricted-primitive-reading}
In \(\mathbf{Set}\),
\[
(x,y)\models \langle R/O/T\rangle\varphi
\]
if and only if there exists an admissible input arrow
\[
a:u\to p_P(x)
\]
such that
\[
R(a),
\qquad
O(\omega_P(a,x)),
\qquad
T(\omega_P(a,x),y),
\]
and
\[
\bigl(
\delta_P(a,x),
\sigma(\omega_P(a,x),y)
\bigr)
\models\varphi.
\]
\end{proposition}

\begin{proof}
\leavevmode
\begin{enumerate}
\item \textbf{Expand the restricted apex.}
The object \((R/O/T)^\sharp\) is the subobject of primitive executions
satisfying the three pulled-back predicates
\[
R^\sharp,\qquad O^\sharp,\qquad T^\sharp.
\]

\item \textbf{Expand the diamond.}
The diamond asserts existence of a restricted primitive execution whose source
is \((x,y)\) and whose target satisfies \(\varphi\).

\item \textbf{Read this in elements.}
Such a restricted primitive execution is exactly an admissible input arrow
\(a:u\to p_P(x)\) satisfying the displayed input, output, transition, and
target conditions.
\end{enumerate}
\end{proof}

\section{Finite choice of subprograms}
\label{sec:relative-finite-choice}

Finite choice is the categorical counterpart of nondeterministic choice in
the non-iterative dynamic fragment \cite{HarelKozenTiurynDynamicLogic}. In the
regular/coherent setting it is interpreted by finite joins of subobjects.

\begin{definition}[Coherent category]
\label{def:relative-coherent-category}
A \textbf{coherent category} is a regular category \(\mathcal E\) such that:
\begin{enumerate}
\item each \(\mathrm{Sub}(X)\) has finite joins;
\item each pullback functor \(f^\ast\) preserves finite joins.
\end{enumerate}
\end{definition}

Assume \(\mathcal E\) is coherent. Let
\[
U,V\hookrightarrow M_{\Phi,Y}
\]
be subprograms of the relative primitive execution span. Their finite choice
is the join
\[
U\vee V\hookrightarrow M_{\Phi,Y}.
\]
For a subobject \(i_U:U\hookrightarrow M_{\Phi,Y}\), write
\[
\langle U\rangle\varphi
:=
\exists_{s i_U}\bigl((t i_U)^\ast\varphi\bigr).
\]
Equivalently, regarding \(U\) as a predicate on \(M_{\Phi,Y}\),
\[
\langle U\rangle\varphi
=
\exists_s\bigl(U\wedge t^\ast\varphi\bigr).
\]

\begin{proposition}[Finite choice law]
\label{prop:relative-finite-choice-law}
For subprograms
\[
U,V\hookrightarrow M_{\Phi,Y},
\]
one has
\[
\langle U\vee V\rangle\varphi
=
\langle U\rangle\varphi
\vee
\langle V\rangle\varphi.
\]
\end{proposition}

\begin{proof}
\leavevmode
\begin{enumerate}
\item \textbf{Regard subprograms as predicates.}
Treat \(U\) and \(V\) as predicates on \(M_{\Phi,Y}\). Then
\[
\langle U\vee V\rangle\varphi
=
\exists_s\bigl((U\vee V)\wedge t^\ast\varphi\bigr).
\]

\item \textbf{Distribute meet over finite join.}
In a coherent category,
\[
(U\vee V)\wedge t^\ast\varphi
=
(U\wedge t^\ast\varphi)
\vee
(V\wedge t^\ast\varphi).
\]

\item \textbf{Use preservation of joins by \(\exists_s\).}
Since \(\exists_s\) is a left adjoint, it preserves finite joins. Hence
\[
\langle U\vee V\rangle\varphi
=
\langle U\rangle\varphi
\vee
\langle V\rangle\varphi.
\]
\end{enumerate}
\end{proof}

\section{Boxes and intuitionistic tests}
\label{sec:relative-boxes}

The box modality is the universal predicate transformer dual to the diamond.
In the categorical setting below, its existence is supplied by the right
adjoints to pullback in first-order Heyting categories
\cite{Jacobs1999CategoricalLogic}. The resulting test law
\[
[\theta?]\varphi=\theta\Rightarrow\varphi
\]
is the intuitionistic form of the usual dynamic-logic test law
\cite{HarelKozenTiurynDynamicLogic}.

\subsection{First-order Heyting structure}
\label{subsec:relative-first-order-heyting}

\begin{definition}[First-order Heyting category]
\label{def:relative-first-order-heyting}
A \textbf{first-order Heyting category} is a coherent category
\(\mathcal E\) such that, for every morphism
\[
f:X\to Y,
\]
the pullback functor
\[
f^\ast:\mathrm{Sub}(Y)\to\mathrm{Sub}(X)
\]
has a right adjoint
\[
\forall_f:\mathrm{Sub}(X)\to\mathrm{Sub}(Y),
\]
and these right adjoints satisfy Beck--Chevalley for pullback squares.
\end{definition}

Assume \(\mathcal E\) is a first-order Heyting category.

\subsection{Box modalities}
\label{subsec:relative-box-modalities}

\begin{definition}[Box modality of a span]
\label{def:relative-span-box}
For a span
\[
R=
(X\xleftarrow{\ell}R\xrightarrow{r}Y),
\]
define
\[
[R]:\mathrm{Sub}(Y)\to\mathrm{Sub}(X)
\]
by
\[
[R]\varphi:=\forall_\ell(r^\ast\varphi).
\]
\end{definition}

\begin{proposition}[Functoriality of boxes]
\label{prop:relative-box-functoriality}
For spans
\[
R:X\nrightarrow Y
\qquad\text{and}\qquad
S:Y\nrightarrow Z,
\]
one has
\[
[1_X]=\mathrm{id}_{\mathrm{Sub}(X)}
\]
and
\[
[R;S]=[R]\circ[S].
\]
\end{proposition}

\begin{proof}
\leavevmode
\begin{enumerate}
\item \textbf{Check the identity span.}
For the identity span, the formula gives
\[
\forall_{1_X}1_X^\ast=\mathrm{id}_{\mathrm{Sub}(X)}.
\]

\item \textbf{Use the composite pullback.}
For
\[
R=(X\xleftarrow{\ell}R\xrightarrow{r}Y),
\qquad
S=(Y\xleftarrow{m}S\xrightarrow{n}Z),
\]
the composite has apex \(R\times_Y S\).

\item \textbf{Apply Beck--Chevalley for \(\forall\).}
The Beck--Chevalley law for the pullback square defining \(R\times_Y S\),
together with functoriality of right adjoints, gives
\[
[R;S]\varphi
=
\forall_\ell r^\ast\forall_m n^\ast\varphi
=
[R]([S]\varphi).
\]
\end{enumerate}
\end{proof}

\subsection{Boxes for tests}
\label{subsec:relative-boxes-tests}

\begin{proposition}[Box law for tests]
\label{prop:relative-test-box}
For a test
\[
\theta?=(S\xleftarrow{\theta}\Theta\xrightarrow{\theta}S),
\]
one has
\[
[\theta?]\varphi
=
\theta\Rightarrow\varphi.
\]
\end{proposition}

\begin{proof}
\leavevmode
\begin{enumerate}
\item \textbf{Expand the box.}
By definition,
\[
[\theta?]\varphi
=
\forall_\theta(\theta^\ast\varphi).
\]

\item \textbf{Test against an arbitrary predicate.}
For any predicate \(\psi\hookrightarrow S\),
\[
\psi\le \forall_\theta(\theta^\ast\varphi)
\]
if and only if
\[
\theta^\ast\psi\le \theta^\ast\varphi.
\]

\item \textbf{Use monicity of \(\theta\).}
Since \(\theta\) is monic, this is equivalent to
\[
\psi\wedge\theta\le\varphi.
\]

\item \textbf{Use the Heyting implication.}
The last inequality is equivalent to
\[
\psi\le\theta\Rightarrow\varphi.
\]
Hence
\[
[\theta?]\varphi=\theta\Rightarrow\varphi.
\]
\end{enumerate}
\end{proof}

\section{Stability under downstream change and Mealy simulation}
\label{sec:relative-modal-stability}

\subsection{Change of downstream action}
\label{subsec:relative-modal-stability-downstream}

Let
\[
f:Y\to Z
\]
be a morphism of right \(\mathbb B\)-actions. It induces
\[
h_f:S_{\Phi,Y}\to S_{\Phi,Z},
\qquad
h_f(x,y)=(x,f(y)),
\]
and
\[
k_f:M_{\Phi,Y}\to M_{\Phi,Z},
\qquad
k_f(a,x,y)=(a,x,f(y)).
\]
The source square is a pullback:
\[
\begin{tikzcd}
M_{\Phi,Y} \arrow[r,"k_f"] \arrow[d,"s_Y"']
&
M_{\Phi,Z} \arrow[d,"s_Z"]
\\
S_{\Phi,Y} \arrow[r,"h_f"']
&
S_{\Phi,Z}
\arrow[ul,phantom,very near start,"\lrcorner"]
\end{tikzcd}
.
\]
Indeed, both \(M_{\Phi,Y}\) and \(M_{\Phi,Z}\) are obtained by pulling back
\(t_A:A_1\to A_0\) along the respective structure maps to \(A_0\), and
\(h_f\) preserves those structure maps.

\begin{proposition}[Diamond stability under downstream change]
\label{prop:relative-diamond-stability-downstream}
Let \(\mathcal E\) be regular. For every subprogram
\[
V\hookrightarrow M_{\Phi,Z}
\]
and every predicate
\[
\psi\hookrightarrow S_{\Phi,Z},
\]
one has
\[
h_f^\ast\langle V\rangle\psi
=
\langle k_f^\ast V\rangle h_f^\ast\psi.
\]
\end{proposition}

\begin{proof}
Let
\[
i:V\hookrightarrow M_{\Phi,Z}
\]
be the subprogram inclusion, and form the pullback
\[
\begin{tikzcd}
k_f^\ast V \arrow[r,"\bar k_f"] \arrow[d,"\bar i"']
&
V \arrow[d,"i"]
\\
M_{\Phi,Y} \arrow[r,"k_f"']
&
M_{\Phi,Z}
\arrow[ul,phantom,very near start,"\lrcorner"]
\end{tikzcd}
.
\]
Because the source square
\[
\begin{tikzcd}
M_{\Phi,Y} \arrow[r,"k_f"] \arrow[d,"s_Y"']
&
M_{\Phi,Z} \arrow[d,"s_Z"]
\\
S_{\Phi,Y} \arrow[r,"h_f"']
&
S_{\Phi,Z}
\end{tikzcd}
\]
is a pullback, the induced square
\[
\begin{tikzcd}
k_f^\ast V \arrow[r,"\bar k_f"] \arrow[d,"s_Y\bar i"']
&
V \arrow[d,"s_Z i"]
\\
S_{\Phi,Y} \arrow[r,"h_f"']
&
S_{\Phi,Z}
\end{tikzcd}
\]
is also a pullback. Moreover target compatibility gives
\[
t_Z i\bar k_f=h_f t_Y\bar i.
\]
Therefore Beck--Chevalley for regular images gives
\[
h_f^\ast\exists_{s_Z i}\bigl((t_Z i)^\ast\psi\bigr)
=
\exists_{s_Y\bar i}\bar k_f^\ast\bigl((t_Z i)^\ast\psi\bigr).
\]
Using target compatibility, the right-hand side becomes
\[
\exists_{s_Y\bar i}\bigl((t_Y\bar i)^\ast h_f^\ast\psi\bigr)
=
\langle k_f^\ast V\rangle h_f^\ast\psi.
\]
\end{proof}

\begin{proposition}[Box stability under downstream change]
\label{prop:relative-box-stability-downstream}
Let \(\mathcal E\) be first-order Heyting. For every subprogram
\[
V\hookrightarrow M_{\Phi,Z}
\]
and every predicate
\[
\psi\hookrightarrow S_{\Phi,Z},
\]
one has
\[
h_f^\ast[V]\psi
=
[k_f^\ast V]h_f^\ast\psi.
\]
\end{proposition}

\begin{proof}
With notation as in the proof of
Proposition~\ref{prop:relative-diamond-stability-downstream}, the square
\[
\begin{tikzcd}
k_f^\ast V \arrow[r,"\bar k_f"] \arrow[d,"s_Y\bar i"']
&
V \arrow[d,"s_Z i"]
\\
S_{\Phi,Y} \arrow[r,"h_f"']
&
S_{\Phi,Z}
\end{tikzcd}
\]
is a pullback, and target compatibility gives
\[
t_Z i\bar k_f=h_f t_Y\bar i.
\]
Hence Beck--Chevalley for the right adjoints \(\forall\) gives
\[
h_f^\ast\forall_{s_Z i}\bigl((t_Z i)^\ast\psi\bigr)
=
\forall_{s_Y\bar i}\bar k_f^\ast\bigl((t_Z i)^\ast\psi\bigr).
\]
Using target compatibility, this is
\[
\forall_{s_Y\bar i}\bigl((t_Y\bar i)^\ast h_f^\ast\psi\bigr)
=
[k_f^\ast V]h_f^\ast\psi.
\]
\end{proof}

\subsection{Change along a Mealy simulation}
\label{subsec:relative-modal-stability-simulation}

Let
\[
\Theta:\Phi\Rightarrow\Psi
\]
be a Mealy simulation, and let \(Y\) be a right \(\mathbb B\)-action. The
induced functor
\[
H_{\Theta,Y}:\mathsf{Prog}_{\Psi,Y}\to\mathsf{Prog}_{\Phi,Y}
\]
has object map
\[
h_{\Theta,Y}:S_{\Psi,Y}\to S_{\Phi,Y}
\]
and arrow map
\[
k_{\Theta,Y}:M_{\Psi,Y}\to M_{\Phi,Y},
\qquad
k_{\Theta,Y}(a,y,z)=
(a,\theta y,\sigma(\alpha_y,z)).
\]
The source square is a pullback:
\[
\begin{tikzcd}
M_{\Psi,Y} \arrow[r,"k_{\Theta,Y}"] \arrow[d,"s_\Psi"']
&
M_{\Phi,Y} \arrow[d,"s_\Phi"]
\\
S_{\Psi,Y} \arrow[r,"h_{\Theta,Y}"']
&
S_{\Phi,Y}
\arrow[ul,phantom,very near start,"\lrcorner"]
\end{tikzcd}
.
\]
Indeed, both \(M_{\Psi,Y}\) and \(M_{\Phi,Y}\) are obtained by pulling back
\(t_A:A_1\to A_0\) along the respective structure maps to \(A_0\), and
\(h_{\Theta,Y}\) preserves those structure maps.

\begin{proposition}[Diamond stability under Mealy simulation]
\label{prop:relative-diamond-stability-simulation}
Let \(\mathcal E\) be regular. For every subprogram
\[
V\hookrightarrow M_{\Phi,Y}
\]
and every predicate
\[
\psi\hookrightarrow S_{\Phi,Y},
\]
one has
\[
h_{\Theta,Y}^\ast\langle V\rangle\psi
=
\langle k_{\Theta,Y}^\ast V\rangle h_{\Theta,Y}^\ast\psi.
\]
\end{proposition}

\begin{proof}
Let
\[
i:V\hookrightarrow M_{\Phi,Y}
\]
be the subprogram inclusion, and form the pullback
\[
\begin{tikzcd}
k_{\Theta,Y}^\ast V \arrow[r,"\bar k_{\Theta,Y}"] \arrow[d,"\bar i"']
&
V \arrow[d,"i"]
\\
M_{\Psi,Y} \arrow[r,"k_{\Theta,Y}"']
&
M_{\Phi,Y}
\arrow[ul,phantom,very near start,"\lrcorner"]
\end{tikzcd}
.
\]
Because the source square
\[
\begin{tikzcd}
M_{\Psi,Y} \arrow[r,"k_{\Theta,Y}"] \arrow[d,"s_\Psi"']
&
M_{\Phi,Y} \arrow[d,"s_\Phi"]
\\
S_{\Psi,Y} \arrow[r,"h_{\Theta,Y}"']
&
S_{\Phi,Y}
\end{tikzcd}
\]
is a pullback, the induced square
\[
\begin{tikzcd}
k_{\Theta,Y}^\ast V
\arrow[r,"\bar k_{\Theta,Y}"]
\arrow[d,"s_\Psi\bar i"']
&
V
\arrow[d,"s_\Phi i"]
\\
S_{\Psi,Y}
\arrow[r,"h_{\Theta,Y}"']
&
S_{\Phi,Y}
\end{tikzcd}
\]
is also a pullback. Since \(H_{\Theta,Y}\) is an internal functor, target
compatibility gives
\[
t_\Phi i\bar k_{\Theta,Y}=h_{\Theta,Y}t_\Psi\bar i.
\]
Therefore Beck--Chevalley for regular images gives
\[
h_{\Theta,Y}^\ast
\exists_{s_\Phi i}\bigl((t_\Phi i)^\ast\psi\bigr)
=
\exists_{s_\Psi\bar i}
\bar k_{\Theta,Y}^\ast\bigl((t_\Phi i)^\ast\psi\bigr).
\]
Using target compatibility, the right-hand side is
\[
\exists_{s_\Psi\bar i}
\bigl((t_\Psi\bar i)^\ast h_{\Theta,Y}^\ast\psi\bigr)
=
\langle k_{\Theta,Y}^\ast V\rangle h_{\Theta,Y}^\ast\psi.
\]
\end{proof}

\begin{proposition}[Box stability under Mealy simulation]
\label{prop:relative-box-stability-simulation}
Let \(\mathcal E\) be first-order Heyting. For every subprogram
\[
V\hookrightarrow M_{\Phi,Y}
\]
and every predicate
\[
\psi\hookrightarrow S_{\Phi,Y},
\]
one has
\[
h_{\Theta,Y}^\ast[V]\psi
=
[k_{\Theta,Y}^\ast V]h_{\Theta,Y}^\ast\psi.
\]
\end{proposition}

\begin{proof}
With notation as in the proof of
Proposition~\ref{prop:relative-diamond-stability-simulation}, the square
\[
\begin{tikzcd}
k_{\Theta,Y}^\ast V
\arrow[r,"\bar k_{\Theta,Y}"]
\arrow[d,"s_\Psi\bar i"']
&
V
\arrow[d,"s_\Phi i"]
\\
S_{\Psi,Y}
\arrow[r,"h_{\Theta,Y}"']
&
S_{\Phi,Y}
\end{tikzcd}
\]
is a pullback, and target compatibility gives
\[
t_\Phi i\bar k_{\Theta,Y}=h_{\Theta,Y}t_\Psi\bar i.
\]
Hence Beck--Chevalley for the right adjoints \(\forall\) gives
\[
h_{\Theta,Y}^\ast
\forall_{s_\Phi i}\bigl((t_\Phi i)^\ast\psi\bigr)
=
\forall_{s_\Psi\bar i}
\bar k_{\Theta,Y}^\ast\bigl((t_\Phi i)^\ast\psi\bigr).
\]
Using target compatibility, this becomes
\[
\forall_{s_\Psi\bar i}
\bigl((t_\Psi\bar i)^\ast h_{\Theta,Y}^\ast\psi\bigr)
=
[k_{\Theta,Y}^\ast V]h_{\Theta,Y}^\ast\psi.
\]
\end{proof}

\section{The one-object specialization}
\label{sec:relative-one-object-specialization}

Let \(M\) and \(N\) be monoid objects in \(\mathcal E\), regarded as
one-object internal categories with the order convention
\[
b\circ a=ba.
\]
A Mealy morphism
\[
\Phi:M\rightsquigarrow N
\]
consists of an object \(P\), a right \(M\)-action
\[
x\cdot m:=\delta_P(m,x),
\]
and an \(N\)-valued output cocycle
\[
\omega_P:M\times P\to N
\]
satisfying
\[
\omega_P(1,x)=1,
\qquad
\omega_P(mn,x)=\omega_P(m,x)\omega_P(n,x\cdot m).
\]

In this subsection we switch from the internal-action notation
\(\sigma(n,y)\) to the usual right-action notation
\[
y\cdot n:=\sigma(n,y).
\]

\begin{proposition}[Classical cascade form]
\label{prop:relative-one-object-cascade}
Let \(Y\) be a right \(N\)-object. Then
\[
\Phi^\ast Y=P\times Y
\]
with right \(M\)-action
\[
(x,y)\cdot m
=
\bigl(x\cdot m,\;y\cdot \omega_P(m,x)\bigr).
\]
\end{proposition}

\begin{proof}
\leavevmode
\begin{enumerate}
\item \textbf{Specialize the coupled state object.}
Since the object objects of \(M\) and \(N\) are terminal, the pullback
\[
P\times_{B_0}Y
\]
is simply \(P\times Y\).

\item \textbf{Specialize the restricted action.}
Definition~\ref{def:relative-restricted-action} gives, in internal-action
notation,
\[
m\cdot_{\mathrm{int}}(x,y)
=
\bigl(\delta_P(m,x),\sigma(\omega_P(m,x),y)\bigr).
\]

\item \textbf{Rewrite in right-action notation.}
With
\[
x\cdot m:=\delta_P(m,x),
\qquad
y\cdot n:=\sigma(n,y),
\]
and
\[
(x,y)\cdot m:=m\cdot_{\mathrm{int}}(x,y),
\]
this becomes
\[
(x,y)\cdot m
=
\bigl(x\cdot m,\;y\cdot \omega_P(m,x)\bigr).
\]
\end{enumerate}
\end{proof}

\begin{proposition}[One-object simulations evaluate by coboundary action]
\label{prop:relative-one-object-simulation-evaluation}
Let
\[
\Theta:\Phi\Rightarrow\Psi
\]
be a one-object Mealy simulation with components
\[
\theta:Q\to P,
\qquad
\alpha:Q\to N.
\]
Then for every right \(N\)-object \(Y\), the evaluated map
\[
Q\times Y\to P\times Y
\]
is
\[
(y,z)\longmapsto(\theta y,z\cdot\alpha_y).
\]
The simulation equation is exactly the equivariance of this map with respect
to the induced right \(M\)-actions. Explicitly, equivariance is the conjunction
of
\[
\theta(y\cdot m)=\theta y\cdot m
\]
and
\[
\alpha_y\,\omega_P(m,\theta y)
=
\omega_Q(m,y)\,\alpha_{y\cdot m}.
\]
\end{proposition}

\begin{proof}
\leavevmode
\begin{enumerate}
\item \textbf{Specialize the evaluated map.}
The general formula from
Definition~\ref{def:relative-evaluated-simulation-map} gives
\[
(y,z)\longmapsto(\theta y,\sigma(\alpha_y,z)).
\]

\item \textbf{Rewrite in right-action notation.}
Since \(z\cdot n:=\sigma(n,z)\), this becomes
\[
(y,z)\longmapsto(\theta y,z\cdot\alpha_y).
\]

\item \textbf{Specialize equivariance.}
Theorem~\ref{thm:relative-simulations-evaluate-equivariantly} says that this
map is equivariant. Written in one-object notation, that equivariance is
exactly the displayed one-object simulation equation.
\end{enumerate}
\end{proof}

\section{The stateless specialization}
\label{sec:relative-stateless-specialization}

Let
\[
F:\mathbb A\to\mathbb B
\]
be an internal functor. Regard \(F\) as the map-representable Mealy morphism
with carrier
\[
A_0\xleftarrow{1_{A_0}}A_0\xrightarrow{F_0}B_0.
\]

\begin{proposition}[Restriction of actions along an internal functor]
\label{prop:relative-stateless-restriction}
Let
\[
Y=(Y\xrightarrow{r}B_0,\sigma)
\]
be a right \(\mathbb B\)-action. Then
\[
F^\ast Y
=
A_0\times_{F_0,B_0,r}Y
\]
with right \(\mathbb A\)-action
\[
a\cdot(v,y)
=
\bigl(s_A(a),\sigma(F_1(a),y)\bigr),
\]
where
\[
t_A(a)=v,
\qquad
r(y)=F_0(v).
\]
\end{proposition}

\begin{proof}
\leavevmode
\begin{enumerate}
\item \textbf{Specialize the carrier.}
For the map-representable carrier, the state component is \(A_0\), with left
leg \(1_{A_0}\) and right leg \(F_0\).

\item \textbf{Specialize the coupled state object.}
The coupled state object becomes
\[
A_0\times_{F_0,B_0,r}Y.
\]

\item \textbf{Specialize the action.}
The state component moves from \(t_A(a)\) to \(s_A(a)\), and the emitted
output arrow is \(F_1(a)\). Hence
\[
a\cdot(v,y)
=
\bigl(s_A(a),\sigma(F_1(a),y)\bigr).
\]
\end{enumerate}
\end{proof}

\begin{corollary}[Terminal stateless case]
\label{cor:relative-terminal-stateless}
For \(Y=\mathbf 1_{\mathbb B}\),
\[
\mathsf{Prog}_{F,\mathbf 1_{\mathbb B}}
\cong
\mathbb A^{\mathrm{op}}.
\]
Under this identification, the input projection is the identity on
\(\mathbb A^{\mathrm{op}}\), and the output labelling is
\[
F^{\mathrm{op}}:\mathbb A^{\mathrm{op}}\to\mathbb B^{\mathrm{op}}.
\]
\end{corollary}

\begin{proof}
\leavevmode
\begin{enumerate}
\item \textbf{Identify the state object.}
The state object is
\[
A_0\times_{F_0,B_0,1_{B_0}}B_0\cong A_0.
\]

\item \textbf{Identify the restricted action.}
The restricted action is the terminal right \(\mathbb A\)-action.

\item \textbf{Identify the action category.}
The action category of the terminal right \(\mathbb A\)-action is
\[
\mathbb A^{\mathrm{op}}.
\]

\item \textbf{Identify the output labelling.}
The output labelling sends an arrow \(a\) to \(F_1(a)\), viewed oppositely.
Thus it is
\[
F^{\mathrm{op}}.
\]
\end{enumerate}
\end{proof}

\begin{proposition}[Natural transformations as evaluated simulations]
\label{prop:relative-stateless-natural-transformations}
Let
\[
\alpha:F\Rightarrow G
\]
be an internal natural transformation, regarded as a Mealy simulation
\[
F\Rightarrow G.
\]
For every right \(\mathbb B\)-action \(Y\), the evaluated map
\[
G^\ast Y\to F^\ast Y
\]
is
\[
(v,y)\longmapsto
(v,\sigma(\alpha_v,y)).
\]
\end{proposition}

\begin{proof}
\leavevmode
\begin{enumerate}
\item \textbf{Identify the state comparison.}
In the stateless case, the state-comparison map is the identity on \(A_0\).

\item \textbf{Identify the output comparison.}
The output-comparison component is the natural transformation
\[
\alpha_v:F_0v\to G_0v.
\]

\item \textbf{Evaluate in the downstream action.}
Evaluating this component in \(Y\) gives
\[
(v,y)\longmapsto
(v,\sigma(\alpha_v,y)).
\]
\end{enumerate}
\end{proof}

\section{Examples}
\label{sec:relative-examples}

The examples below indicate how the abstract construction specializes to
classical Mealy cascades and relational dynamic logic. The Mealy terminology
follows the categorical lineage represented by Paré's Mealy morphisms
\cite{PareMealyMorphisms}, while the relational dynamic-logic reading follows
the standard program-modality interpretation
\cite{HarelKozenTiurynDynamicLogic}.

\begin{example}[Classical Mealy cascade]
\label{ex:relative-classical-mealy-cascade}
Let \(\mathcal E=\mathbf{Set}\), and suppose \(\mathbb A\) and \(\mathbb B\)
are one-object internal categories corresponding to monoids \(M\) and \(N\).
A Mealy morphism consists of a set \(P\), a right \(M\)-action, and an
\(N\)-valued output cocycle
\[
\omega:M\times P\to N.
\]
A downstream right \(N\)-action \(Y\) produces the coupled state set
\[
P\times Y
\]
with right \(M\)-action
\[
(x,y)\cdot m
=
\bigl(x\cdot m,\;y\cdot\omega(m,x)\bigr).
\]
This is the cascade of a Mealy transducer followed by an \(N\)-system.
\end{example}

\begin{example}[Simulation as state-dependent output correction]
\label{ex:relative-simulation-output-correction}
In the one-object group-valued case, let \(N=G\) be a group. A simulation
between two \(G\)-valued output cocycles has a state map
\[
\theta:Q\to P
\]
and a correction term
\[
\alpha:Q\to G.
\]
The simulation equation says that the corrected \(P\)-output and the corrected
\(Q\)-output agree:
\[
\alpha_y\,\omega_P(m,\theta y)
=
\omega_Q(m,y)\,\alpha_{y\cdot m}.
\]
For every right \(G\)-set \(Y\), this correction is evaluated by
\[
(y,z)\longmapsto(\theta y,z\cdot\alpha_y).
\]
Thus output-comparison data becomes an actual transition in each downstream
context.
\end{example}

\begin{example}[Relational semantics of the non-iterative dynamic fragment]
\label{ex:relative-relational-pdl}
Let \(\mathcal E=\mathbf{Set}\). This is the ordinary relational semantics of
the non-iterative fragment of dynamic logic
\cite{HarelKozenTiurynDynamicLogic}. A span
\[
S\xleftarrow{\ell}R\xrightarrow{r}S
\]
determines the binary relation
\[
x\,\overline R\,y
\]
if and only if there exists \(\rho\in R\) such that
\[
\ell(\rho)=x,
\qquad
r(\rho)=y.
\]
The diamond becomes the relational dynamic-logic modality:
\[
x\models\langle R\rangle\varphi
\]
if and only if there exists \(y\) such that
\[
x\,\overline R\,y
\]
and
\[
y\models\varphi.
\]
In the relative Mealy setting, this relational semantics is applied to the
coupled state space
\[
S_{\Phi,Y}=P\times_{B_0}Y.
\]
\end{example}

\section{Chapter summary}
\label{sec:relative-summary}

\begin{theorem}[Relative Mealy dynamic span logic]
\label{thm:relative-mealy-dynamic-span-logic}
The constructions of this chapter form the following hierarchy.

\begin{enumerate}
\item Right internal \(\mathbb A\)-actions are Eilenberg--Moore algebras for
the monad induced by \(\mathbb A\) on \(\mathcal E/A_0\).

\item A right action \(X\) has an action category
\[
\int_{\mathbb A}X
\]
and a canonical input-discrete projection
\[
\int_{\mathbb A}X\to\mathbb A^{\mathrm{op}}.
\]

\item Every downstream right \(\mathbb B\)-action \(Y\) evaluates
output-comparison Kleisli cells
\[
Q\to B\circ P
\]
as maps
\[
Q\times_{B_0}Y\to P\times_{B_0}Y.
\]

\item A Mealy morphism
\[
\Phi:\mathbb A\rightsquigarrow\mathbb B
\]
restricts every downstream action \(Y\) to a right \(\mathbb A\)-action
\[
\Phi^\ast Y.
\]

\item A Mealy simulation
\[
\Theta:\Phi\Rightarrow\Psi
\]
evaluates contravariantly to an equivariant map
\[
\Psi^\ast Y\to\Phi^\ast Y.
\]

\item The relative program category is
\[
\mathsf{Prog}_{\Phi,Y}
=
\int_{\mathbb A}\Phi^\ast Y.
\]

\item The primitive execution span
\[
\mathbf m_{\Phi,Y}:S_{\Phi,Y}\nrightarrow S_{\Phi,Y}
\]
is a monoid object in
\[
\operatorname{EndSpan}_{\mathcal E}(S_{\Phi,Y}).
\]

\item The construction
\[
\Phi\longmapsto\mathsf{Prog}_{\Phi,Y}
\]
is contravariant in Mealy simulations.

\item There are structural functors
\[
I_{\Phi,Y}:\mathsf{Prog}_{\Phi,Y}\to\mathbb A^{\mathrm{op}}
\]
and
\[
\Lambda_{\Phi,Y}:\mathsf{Prog}_{\Phi,Y}\to\int_{\mathbb B}Y.
\]

\item A simulation
\[
\Theta:\Phi\Rightarrow\Psi
\]
induces an input-strict functor
\[
H_{\Theta,Y}:
\mathsf{Prog}_{\Psi,Y}\to\mathsf{Prog}_{\Phi,Y}
\]
and a downstream comparison natural transformation
\[
\Lambda_{\Psi,Y}
\Rightarrow
\Lambda_{\Phi,Y}H_{\Theta,Y}.
\]

\item If \(\mathcal E\) is regular, every relative execution span acts on
subobject predicates by the diamond
\[
\langle R\rangle\varphi=\exists_\ell(r^\ast\varphi),
\]
using the regular-image interpretation of existential quantification
\cite{Jacobs1999CategoricalLogic,FongSpivakRegularRelational}.

\item If \(\mathcal E\) is coherent, finite joins of relative subprograms give
finite nondeterministic choice.

\item If \(\mathcal E\) is first-order Heyting, every relative execution span
also acts by the box
\[
[R]\varphi=\forall_\ell(r^\ast\varphi),
\]
and tests satisfy
\[
[\theta?]\varphi=\theta\Rightarrow\varphi.
\]
This is the first-order Heyting-categorical form of the box and test semantics
for the non-iterative dynamic fragment
\cite{Jacobs1999CategoricalLogic,HarelKozenTiurynDynamicLogic}.

\item Modalities are stable under downstream action maps and under Mealy
simulations, after pulling subprograms back along the induced program
functors.

\item The one-object specialization gives the cascade action
\[
(x,y)\cdot m
=
\bigl(x\cdot m,\;y\cdot\omega_P(m,x)\bigr),
\]
and the stateless specialization recovers restriction of actions along
internal functors.
\end{enumerate}
\end{theorem}

\begin{proof}
\leavevmode
\begin{enumerate}
\item \textbf{Actions and action categories.}
Items (1)--(2) are
Proposition~\ref{prop:relative-actions-em},
Theorem~\ref{thm:relative-action-category-internal}, and
Theorem~\ref{thm:relative-actions-input-discrete}.

\item \textbf{Downstream evaluation and restrictions.}
Items (3)--(5) are
Proposition~\ref{prop:relative-evaluation-functorial},
Theorem~\ref{thm:relative-mealy-restriction-actions}, and
Theorem~\ref{thm:relative-simulations-evaluate-equivariantly}.

\item \textbf{Program categories and labelled functoriality.}
Items (6)--(10) are
Theorem~\ref{thm:relative-program-categories},
Theorem~\ref{thm:relative-simulations-program-functors},
Theorem~\ref{thm:relative-input-downstream-functors},
Theorem~\ref{thm:relative-output-comparison-natural}, and
Corollary~\ref{cor:relative-labelled-functor-simulation}.

\item \textbf{Modal structure.}
Items (11)--(14) are the span predicate-transformer constructions of
Sections~\ref{sec:relative-predicate-transformers}--\ref{sec:relative-modal-stability}.

\item \textbf{Specializations.}
Item (15) is
Proposition~\ref{prop:relative-one-object-cascade} together with
Proposition~\ref{prop:relative-stateless-restriction}.
\end{enumerate}
\end{proof}

\begin{remark}[Final conceptual summary]
\label{rem:relative-final-summary}
A Mealy morphism
\[
\Phi:\mathbb A\rightsquigarrow\mathbb B
\]
is a stateful interpretation of \(\mathbb A\)-inputs as
\(\mathbb B\)-outputs, in the same broad categorical lineage as Paré's Mealy
morphisms of enriched categories \cite{PareMealyMorphisms}. A downstream right
\(\mathbb B\)-action \(Y\) gives those outputs a place to act. The coupled
state object
\[
P\times_{B_0}Y
\]
therefore carries a right \(\mathbb A\)-action, and its action category is the
relative program category
\[
\mathsf{Prog}_{\Phi,Y}.
\]
A Mealy simulation carries output-comparison arrows. The downstream action
evaluates those arrows, so simulations become equivariant maps after
evaluation. Dynamic span logic is then the non-iterative modal logic of spans
on the coupled state object, functorial under downstream change and
contravariantly functorial under Mealy simulation.
\end{remark}